\documentclass{../template/note}

\author{Oliver West}
\title{Vectors and Matrices}
\date{}

\begin{document}
    
    \maketitle

    \tableofcontents

    \lecture{Lecture 1 -- 09/10/25}
    
    \section{Complex numbers}

    \subsection{Definition of complex numbers}

    We construct $\CC$ by adding $i$ to $\RR$, where $i^2=-1$. Any $z \in \CC$ then has the form $x + iy$, $x$, $y \in \RR$. The \emph{real part} $\Re(z) = x$, the \emph{imaginary part} $\Im(z) = y$.

    \subsection{Properties of complex numbers}

    $z$, $z_1$ and $z_2$ are complex numbers throughout, given by $x+iy$, $x_1+iy_1$ and $x_2+iy_2$ respectively.

    \begin{enumerate}[label=\protect\circled{\arabic*}]
        \item \emph{Addition.} $z_1 \pm z_2 = (x_1 \pm x_2) + i(y_1 \pm y_2)$.
        \item \emph{Multiplication.} $z_1z_2 = (x_1x_2-y_1y_2)+i(x_1y_2+x_2y_1)$.
    \end{enumerate}
    Note that both of these operations are associative and commutative.
    \begin{enumerate}[label=\protect\circled{\arabic*},resume]
        \item \emph{Identity.} Additive identity 0, multiplicative identity 1 -- $(\CC, +)$ is an abelian group with identity 0.
        \item \emph{Inverse.} For any complex $z \neq 0$, there exists $z^{-1} \in \CC$ s.t. $zz^{-1} = 1$. In particular,
        \begin{align*}
            z^{-1} = \frac{x-iy}{x^2+y^2}
        \end{align*}
        It follows that $(\CC \setminus \{0\}, \cdot)$ is an abelian group with identity 1.
    \end{enumerate}
    Note also that multiplication distributes over addition.
    \begin{enumerate}[label=\protect\circled{\arabic*},resume]
        \item \emph{Conjugation.} The \emph{complex conjugate} of $z$, $\bar z$ or $z^*$, is given by $x-iy$. Observe $\Re(z) = \frac{1}{2}(z+\bar z)$, $\Im(z) = \frac{1}{2i}(z-\bar z)$. Further important identities relating to conjugation:
        \begin{enumerate}
            \item $\bar{\bar z} = z$;
            \item $\overline{z_1 + z_2} = \overline{z_1} + \overline{z_2}$;
            \item $\overline{z_1z_2} = \overline{z_1}\cdot\overline{z_2}$.
        \end{enumerate}
        
        \item \emph{Modulus.} $r = |z| = \sqrt{x^2+y^2}$.
        
        \item \emph{Argument.} The \emph{(principal)} argument $\theta = \arg z$ of a complex number $z \neq 0$ satisfies $z = r(\cos \theta + i \sin \theta)$. Furthermore, $-\pi < \theta \leq \pi$.
        
        We also define the \emph{multivalued argument} $\Arg z = \{\arg(z) + 2\pi n : n \in \ZZ\}$.
        
        Observe also $\cos \theta = \frac{x}{\sqrt{x^2+y^2}}$, $\sin \theta = \frac{y}{\sqrt{x^2+y^2}}$, $\tan \theta = \frac y x$.

    \end{enumerate}

    \begin{remark*}
        Observe
        \begin{itemize}
            \item $\RR \subset \CC$;
            \item A complex number $z$ is \emph{purely imaginary} if $\Re(z) = 0$;
            \item The representation of a complex number in terms of its real and imaginary parts is unique.
        \end{itemize}
    \end{remark*}
    
    \subsubsection{Further properties}

    $(\CC, +, \cdot)$ is an example of a \emph{field}.

    \begin{theorem}[Fundamental theorem of algebra]
        Let $p(z) = c_nz^n + \dots + c_0$, $c_i \in \CC$, $c_n \neq 0$ be a polynomial of degree $n$. Then there exist $\alpha_1, \dots, \alpha_n$ all in $\CC$ such that $p(z) = c_n(z-\alpha_1)\dots(z-\alpha_n)$. In particular, the equation $p(z)=0$ has at least one root in $\CC$, and $n$ roots when counted with multiplicity.
        \label{fta}
    \end{theorem}

    \noindent Some properties of modulus:

    \begin{itemize}
        \item $|z_1z_2| = |z_1||z_2|$;
        \item $|z_1 + z_2| \leq |z_1| + |z_2|$ \emph{(Triangle inequality)};
        \item $|z_1 - z_2| \geq \left||z_1|-|z_2|\right|$.
    \end{itemize}

    \subsection{The complex plane}

    We represent $z = x+iy$ with the point $(x, y)$.

    \begin{figure}[ht]
        \centering
        \incfig{0.4\textwidth}{argand-basic}
        \caption{A complex number $z = x+iy$ in the complex plane.}
        \label{fig:argand-basic}
    \end{figure}
    
    \begin{figure}[ht]
        \centering
        \hfill
        \begin{subfigure}[t]{0.4\textwidth}
            \incfig{\textwidth}{argand-addition}
            \caption{Addition of complex numbers.}
            \label{fig:argand-addition}
        \end{subfigure}
        \hfill
        \begin{subfigure}[t]{0.35\textwidth}
            \incfig{\textwidth}{argand-conjugation}
            \caption{Conjugation of complex numbers.}
            \label{fig:argand-conjugation}
        \end{subfigure}
        \hfill
    \end{figure}

    \lecture{Lecture 2 -- 11/10/25}

    \begin{theorem}[de Moivre]
        \label{de-moivre-theorem}
        $(\cos \theta + i \sin \theta)^n = \cos n\theta + i \sin n\theta$.
    \end{theorem}

    \noindent Established by induction together with the following lemma:

    \begin{lemma}
        If $z_1=r_1(\cos \theta_1 + i \sin \theta_1)$ and $z_2=r_2(\cos \theta_2 + i \sin \theta_2)$, \\ then $z_1z_2 = r_1r_2\left(\cos(\theta_1 + \theta_2) + i \sin(\theta_1+\theta_2)\right)$.
    \end{lemma}

    \begin{proof}[Proof of lemma]
        \begin{align*}
            z_1z_2 &= r_1(\cos \theta_1 + i \sin \theta_1) \cdot r_2(\cos \theta_2 + i \sin \theta_2) \\
            &= r_1r_2\left[(\cos \theta_1 \cos \theta_2 - \sin \theta_1 \sin \theta_2)+i(\cos \theta_1 \sin \theta_2 + \cos \theta_2 \sin \theta_1)\right] \\
            &= r_1r_2(\cos(\theta_1 + \theta_2) + i \sin(\theta_1 + \theta_2))
        \end{align*} 
    \end{proof}

    \begin{proof}[Proof of theorem] We consider cases.
        \begin{itemize}
            \item If $n=0$, $z^0 = 1 = \cos 0 + i \sin 0$, so works (base case).
            \item If $n \in \NN$, we induct.
            Suppose the IH holds for $n$, that is $(\cos \theta + i \sin \theta)^n = \cos n\theta + i \sin n\theta$. Now
            \begin{align*}
                (\cos \theta + i \sin \theta)^{n+1} &= (\cos \theta + i \sin \theta)^n(\cos \theta + i \sin \theta) \\
                &= (\cos n\theta + i \sin n\theta)(\cos \theta + i \sin \theta) \\
                &= \cos [(n+1)\theta] + i \sin [(n+1)\theta]
            \end{align*}
            by the lemma.
            \item If $n < 0$, we write $n = -m$ with $m \in \NN$. Then
            \begin{align*}
                (\cos \theta + i \sin \theta)^n &= (\cos \theta + i \sin \theta)^{-m} \\
                &= [(\cos \theta + i \sin \theta)^m]^{-1} \\
                &= (\cos m\theta + i \sin m\theta)^{-1} \\
                &= \frac{\cos m\theta - i \sin m\theta}{1} \\
                &= \cos n \theta + i \sin n \theta
            \end{align*}
        \end{itemize}
    \end{proof}
    
    \subsection{Exponential and trigonometric functions}
    
    \begin{definition}
        For $z \in \CC$,
        \begin{align*}
            \exp z = e^z = \sum_{n=0}^\infty \frac{z^n}{n!}
        \end{align*}
        This converges for any $z \in \CC$.        
    \end{definition}
    
    Fundamental properties of the exponential are as follows:
    \begin{itemize}
        \item $e^ze^w=e^{z+w}$ for any $z$, $w \in \CC$.
        \item If $z \in \RR$, this reduces to the usual exponential, so in particular $e^0=1$.
        \item $(e^z)^n = e^{nz}$ for any $z \in \CC$ and any $n \in \ZZ$. Note this does \emph{not} necessarily hold for $n \in \CC$. 
    \end{itemize}
    
    \begin{definition} We define our trigonometric functions as follows:
        
        
        \begin{itemize}
            \item We define cosine as follows:
            \begin{align*}
                \cos z &= \frac 1 2 (e^{iz}+e^{-iz}) \\
                &= \frac 1 2 \left(\sum_{n=0}^\infty \frac{(iz)^n}{n!} + \sum_{n=0}^\infty \frac{(-iz)^n}{n!} \right) \\
                &= \frac 1 2 \left(\sum_{n=0}^\infty \frac{(iz)^n}{n!} + \frac{(-iz)^n}{n!} \right) \\
                &= \sum_{n=0}^\infty (-1)^n \frac{z^{2n}}{(2n)!}
            \end{align*}
            \item We define sine as follows:
            \begin{align*}
                \sin z &= \frac 1 {2i} (e^{iz}-e^{-iz}) \\
                &= \frac 1 {2i} \left(\sum_{n=0}^\infty \frac{(iz)^n}{n!} - \frac{(-iz)^n}{n!} \right) \\
                &= \sum_{n=0}^\infty (-1)^n \frac{z^{2n+1}}{(2n+1)!}
            \end{align*}
        \end{itemize}
        
        
    \end{definition}
    
    Once again, for real $z$, these reduce to the normal definition.

    From these definitions the following result is evident:
    
    \begin{theorem}[Euler's formula]
        $e^{iz} = \cos z + i \sin z$ for any $z \in \CC$.
    \end{theorem}

    Note that \cref{de-moivre-theorem} is immediate from this. Also, if $x \in \RR$, $\Re(e^{ix}) = \cos x$ and $\Im(e^{ix}) = \sin x$.

    \begin{corollary}[Euler's identity]
        In particular, taking $z = \pi$, we have $e^{i \pi} = -1$.
    \end{corollary}

    \begin{lemma}
        $e^z = 1 \iff z = 2\pi n i$ for some $n \in \ZZ$.
    \end{lemma}
    
    \begin{proof}
        Write $z = x + iy$, so
        \begin{align*}
            e^z &= e^xe^{iy} \\
            &= e^x(\cos y + i \sin y)
        \end{align*}
        For the forward direction, suppose $e^x(\cos y + i \sin y) = 1$. Clearly the pure imaginary part must be 0, so $e^x \sin y = 0 \implies \sin y = 0$ ($x$ real implies $e^x > 0$). Since $y$ is real, we must have $y = n\pi$ for some $n \in \ZZ$.

        Also, $e^x \cos y = 1$, but $y = n \pi \implies \cos y = \pm 1$. Since $e^x$ is positive for real $x$, we must have $\cos y = 1$, and thus $n$ even, so $n = 2m$ and $y = 2m\pi$ as required. From here we see $e^x = 1$ with $x$ real, so $x = 0$, so $z = 0 + i \cdot 2 m \pi = 2 \pi m i$ as required.

        For the backward direction, simply substitute $z = 2\pi n i$ and it works.
    \end{proof}
    
    \subsection{Roots of unity}

    Let $z = re^{i \theta}$, and suppose $z^N = 1$ for some $N \in \NN$. Then
    \begin{align*}
        z^N &= r^Ne^{iN\theta} = 1 \\
        &= r^N(\cos N\theta + i \sin N\theta)
    \end{align*}

    Then a similar separation of modulus and argument gives $r^n = 1 \implies r = 1$ (since $r \in \RR_{\geq 0}$). Also, by the lemma, $N \theta = 2 \pi n$, $n \in \ZZ$, and so $\theta = 2\pi n / N$. So
    \begin{align*}
        z &= e^{2\pi n i / N} && n = 0, \dots, N-1 \\
        &= (e^{2\pi i / N})^n \\
        &= \omega^n && \omega = e^{2 \pi i / N}
    \end{align*}
    
    These are the $N^\text{th}$ roots of unity -- there are $N$ of them.

    \subsection{Logarithms and complex powers}

    \begin{definition}[Complex logarithm]
        Let $z \in \CC$, $z \neq 0$, then the \emph{principal logarithm} $\log z = \log|z| + i \arg{z}$. The \emph{multivalued logarithm} $\Log z = \{\log z + 2\pi n i : n \in \ZZ\}$ -- note the analogy to the principal and multivalued arguments.
    \end{definition}
    
    \begin{definition}[Complex exponentiation]
        Let $z$, $\alpha \in \CC$, then $z^\alpha = e^{\alpha \log z}$. This is the principal value -- taking other values of $\log$ can give a different result.
    \end{definition}
    
    \begin{example} Some examples.
        \begin{itemize}
            \item $i = e^{i\pi/2}$. Thus $\arg i = \pi / 2$, $\mathrm{Arg} i = \{\pi/2 + 2\pi n : n \in \ZZ\}$. So $\log i = \log (1 \cdot e^{i\pi/2}) = \frac{i\pi}{2}$. This is the principal logarithm -- adding multiples of $2 \pi i$ gives the other logarithms.
            \item $1 + i = \sqrt 2 e^{i \pi /4}$. So $\log(1 + i) = \log(\sqrt 2) + \frac{i \pi}{4}$ -- again, we can add multiples.
            \item $\sqrt{1+i} = (\sqrt 2 e^{i \pi /4})^{1/2} = e^{\frac{1}{2}\left(\log(\sqrt 2) + \frac{i \pi}{4}+2n\pi i\right)} = \sqrt[4]{2} e^{i \left(\pi/8 + n \pi\right)}$. The principal value is of course given by $n = 0$.
        \end{itemize}
        
        
    \end{example}
    
    \subsection{Lines and circles}

    \subsubsection{Lines}

    Given a point $z_0$ and a direction $\omega$ in $\CC$, a line in the complex plane is given by $z = z_0 + \lambda \omega$ for $\lambda \in \RR$. Then
    \begin{align*}
        z = z_0 + \lambda \omega \implies \overline \omega z &= \overline \omega z_0 + \lambda \omega \overline \omega \\
        \overline z = \overline{z_0} + \lambda \overline \omega \implies \omega \overline z &= \omega \overline{z_0} + \lambda \omega \overline \omega \\[-10pt]
        \cline{1-2}
        \Aboxed{\overline{\omega}z-\omega\overline{z} &= \overline{\omega}z_0-\omega\overline{z_0}}
    \end{align*}

    \subsubsection{Circles}

    Given a centre $c \in \CC$ and radius $\rho > 0$, $\rho \in \RR$, a circle is given by $z = c + \rho e^{i \theta}$, $\theta \in \RR$. Equivalently, 
    \begin{align*}
        \rho&=|z-c|\\
        \iff\rho^2&=|z-c|^2\\
        &= (z-c)(\overline{z}-\overline{c})\\
        &= z\overline{z}-z\overline{c}-\overline{z}c-c\overline{c} \\
        &= |z|^2-z\overline{c}-\overline{z}c-|c|^2\\
        \iff\rho^2-|c|^2&=|z|^2-z\overline{c}-\overline{z}c
    \end{align*}

    \lecture{Lecture 3 -- 14/10/25}
    
    \section{Vectors}

    \begin{definition}
        A \emph{vector} is specified by a positive magnitude and a direction in space.

        We can represent a vector as a line segment between 2 points $A$ and $B$, $\vec v = \overrightarrow{AB}$. Then $\vec v$ has length $|\vec v|$ and direction from $A$ to $B$. If we choose $O$ as the origin, then a point $A$ has position vector $\vec a = \overrightarrow{OA}$.
    \end{definition}

    \subsection{Vector spaces}
    
    \begin{definition}
        A \emph{vector space} $V$ over $\RR$ or $\CC$ is a collection of vectors $\vec v \in V$ together with two operations:
        \begin{itemize}
            \item \emph{addition} of two vectors;
            \item \emph{multiplication} by a scalar from $\RR$ or $\CC$ respectively.
        \end{itemize}
        defined as follows
        \begin{enumerate}[label=\protect\circled{\arabic*}]
            \item \emph{Vector addition}. Suppose $\vec a$ and $\vec b$ are the position vectors of $A$ and $B$ respectively. Construct the parallelogram $OACB$, then $\vec a + \vec b = \vec c$, where $\vec c$ is the position vector of $C$.
            \begin{figure}[ht]
                \centering
                \incfig{0.4\textwidth}{vector-addition}
                \caption{Vector addition}
                \label{fig:vector-addition}
            \end{figure}
            \begin{itemize}
                \item \emph{Commutativity}: $\vec a + \vec b = \vec b + \vec a$
                \item \emph{Associativity}: $(\vec a + \vec b) + \vec c = \vec a + (\vec b + \vec c)$
                \item \emph{Identity}: There is a vector $\vec 0$ such that for any vector $\vec a, \vec a + \vec 0 = \vec a$.
                \item \emph{Inverse}: For any vector $\vec a$, there exists a vector $-\vec a$ such that $\vec a + -\vec a = \vec 0$.
            \end{itemize}
            More concisely, vectors under addition form an abelian group.

            \item \emph{Scalar multiplication}. Given $\vec a$ the position vector of a point $A$, and a scalar $\lambda \in \RR$, $\lambda \vec a$ is the position vector of a point on the line $OA$ with distance from the origin given by $|\lambda||\vec a|$ and direction depending on the sign of $\lambda$.
            
            \begin{figure}[ht]
                \centering
                \incfig{0.6\textwidth}{vector-multiplication}
                \caption{Vector multiplication}
                \label{fig:vector-multiplication}
            \end{figure}
            
            \begin{itemize}
                \item $\lambda(\vec a + \vec b) = \lambda \vec a + \lambda \vec b$
                \item $(\lambda + \mu)\vec a = \lambda \vec a + \mu \vec a$
                \item $\lambda(\mu \vec a) = (\lambda \mu)\vec a$
                \item $1 \cdot \vec a = \vec a$.
            \end{itemize}
            
            
            
            
            
        \end{enumerate}
    \end{definition}
    
    \begin{definition}
        A \emph{unit vector} is a vector with length 1. We denote it $\hvec v$.
    \end{definition}
    
    \begin{definition}
        Let $\vec a$, $\vec b$ be vectors and $\alpha$, $\beta \in \RR$ be scalars. Then $\alpha \vec a + \beta \vec b$ is a \emph{linear combination} of $\vec a$ and $\vec b$. The set of all linear combinations is the \emph{span} of $\vec a$ and $\vec b$, given by $\mathrm{span}\{\vec a, \vec b\} = \{\alpha \vec a + \beta \vec b: \alpha, \beta \in \RR\}$.
    \end{definition}
    
    \begin{definition}
        We say that $\vec a$ and $\vec b$ are \emph{parallel}, $\vec a \parallel \vec b$, if $\vec a = \lambda \vec b$ or $\vec b = \lambda \vec a$ for some $\lambda \in \RR$.
    \end{definition}
    
    \begin{remark}
        $\lambda$ may be 0, so $\vec 0 \parallel \vec a$ for any $\vec a$.
    \end{remark}
    
    If $\vec a \nparallel \vec b$, then $\mathrm{span}\{\vec a, \vec b\}$ is a plane through $OAB$.

    \begin{example}
        $\RR^n$ is a vector space with addition (component by component) and scalar multiplication.

        We can represent $\RR$ by a line, and we might be inclined to conclude that any line is a vector space. But this may not be true, as the line must contain $\vec 0$. In particular, the line must pass through the origin.
    \end{example}
    
    \subsection{The scalar product}

    Let us first consider the usual scalar product in $\RR^n$. Let $\vec a$, $\vec b \in \RR^n$ and $\theta$ the angle between them. Then the scalar product is given by $\vec a \cdot \vec b = |\vec a||\vec b| \cos \theta$.

    \begin{figure}[ht]
        \centering
        \incfig{0.4\textwidth}{dot-product}
        \caption{The scalar product -- first take the length of $\vec b$ projected onto $\vec a$, then scale up by $|\vec a|$.}
        \label{fig:dot-product}
    \end{figure}
    
    If $\vec a$ or $\vec b = \vec 0$, $\theta$ is not defined, but we say $\vec a \cdot \vec b = 0$. Intuitively, it is the product of the `parallel parts' of $\vec a$ and $\vec b$.

    Some properties that hold for real vector spaces:

    \begin{itemize}
        \item $\vec a \cdot \vec b = \vec b \cdot \vec a$;
        \item $\vec a \cdot \vec a = |\vec a|^2 \geq 0$, with equality if and only if $\vec a = \vec 0$.
        \item $(\lambda \vec a) \cdot \vec b = \lambda (\vec a \cdot \vec b) = \vec a \cdot (\lambda \vec b)$.
        \item $\vec a \cdot (\vec b + \vec c) = \vec a \cdot \vec b + \vec a \cdot \vec c$
    \end{itemize}
    
    We say that $\vec a$, $\vec b$ are \emph{orthogonal} or \emph{perpendicular}, denoted by $\vec a \perp \vec b$, if $\vec a \cdot \vec b = 0$.

    It follows from this that $\vec 0$ is orthogonal to any other vector.

    We can also write the projection of $\vec b$ onto $\vec a$ as $(|\vec b| \cos \theta) \hvec a = (\hvec a \cdot \vec b) \hvec a$.
    
    The cosine rule follows as (ADD DIAGRAM)
    
    \begin{align*}
        |\overrightarrow{BC}|^2 &= |\overrightarrow{AB} - \overrightarrow{AC}|^2 \\
        &= (\overrightarrow{AB} - \overrightarrow{AC}) \cdot (\overrightarrow{AB} - \overrightarrow{AC}) \\
        &= |\overrightarrow{AB}|^2 + |\overrightarrow{AC}|^2 - 2\overrightarrow{AB} \cdot \overrightarrow{AC} \\
        &= |\overrightarrow{AB}|^2 + |\overrightarrow{AC}|^2 - 2|\overrightarrow{AB}|| \overrightarrow{AC}|\cos\theta \\
    \end{align*}
    
    This definition extends to any other real vector space as long as the axioms are satisfied.

    \begin{definition}
        For a real vector space $V$, an \emph{inner product} or \emph{scalar product} is a map $\cdot: V^2 \to \RR$ satisfying:
        \begin{itemize}
            \item \emph{Symmetry}: $\vec a \cdot \vec b = \vec b \cdot \vec a$;
            \item \emph{Linearity in second argument}: $\vec a \cdot (\lambda \vec b + \mu \vec c) = \lambda \vec a \cdot \vec b + \mu \vec a \cdot \vec c$;
            \item \emph{Positive-definiteness}: $\vec a \cdot \vec a \geq 0$, with equality if and only if $\vec a = \vec 0$.
        \end{itemize}
        Can also be denoted by $\langle \vec x | \vec y \rangle$.
        
    \end{definition}
    
    \begin{definition}
        The \emph{norm} of a vector $\vec a$, written $|\vec a|$ or $||\vec a||$, is given by $\sqrt{\vec a \cdot \vec a}$.
    \end{definition}

    \begin{theorem}[Cauchy-Schwarz inequality]
        Let $V$ be a real vector space with an inner product $\cdot$ which induces a norm $|\cdot|$. For all $\vec x$, $\vec y \in V$,
        \[
            |\vec x \cdot \vec y| \leq |\vec x|\cdot|\vec y|
        \]
        Furthermore, equality holds if and only if $\vec x \parallel \vec y$.
    \end{theorem}
    
    \begin{proof}
        Consider the expression $|\vec x - \lambda \vec y|^2$ for some real parameter $\lambda$. Then
        \begin{align*}
            0 &\leq |\vec x - \lambda \vec y|^2 \\
            &= (\vec x - \lambda \vec y) \cdot (\vec x - \lambda \vec y) \\
            &= |\vec x|^2 + \lambda^2 |\vec y|^2 - 2 \lambda \vec x \cdot \vec y
        \end{align*}
        This is a quadratic inequality in $\lambda$. In order for the quantity to always be $\geq 0$, it must never cross the $x$-axis, so it can have at most one real root, so the discriminant is $\leq 0$. Thus
        \begin{align*}
            4(\vec x \cdot \vec y)^2 - 4 |\vec x|^2 |\vec y|^2 &\leq 0 \\
            \implies (\vec x \cdot \vec y)^2 &\leq |\vec x|^2 |\vec y|^2 \\
            \implies |\vec x \cdot \vec y| &\leq |\vec x|\cdot|\vec y|
        \end{align*}
        Observe that equality holds at the start only if $\vec x = \lambda \vec y$, giving the equality condition.
    \end{proof}
    
    \begin{remark}
        One may wish to prove this result with the fact that cosine is always less than 1 in magnitude. However, this does not generalise to an arbitrary real vector space.
    \end{remark}
    
    \begin{remark}
        We are now justified in saying $\vec x \cdot \vec y = |\vec x||\vec y| \cos \theta$, since Cauchy-Schwarz ensures $-1 \leq \cos \theta \leq 1$, and we can use this to define the angle between $\vec x$ and $\vec y$ in $\RR^n$
    \end{remark}
    
    \begin{corollary}[Triangle inequality]
        For vectors $\vec x$, $\vec y$ in a real vector space with inner product $\cdot$ and induced norm $|\cdot|$, $|\vec x + \vec y| \leq |\vec x| + |\vec y|$.
    \end{corollary}
    
    \begin{proof}
        \begin{align*}
            |\vec x + \vec y|^2 &= (\vec x + \vec y) \cdot (\vec x + \vec y) \\
            &= |\vec x|^2 + |\vec y|^2 + 2 \vec x \cdot \vec y \\
            &\leq |\vec x|^2 + |\vec y|^2 + 2 |\vec x| \cdot |\vec y| && \text{(Cauchy-Schwarz)}\\
            &= (|\vec x| + |\vec y|)^2
        \end{align*}
        Since the norm is always $\geq 0$, we can take the square root to get $|\vec x + \vec y| \leq |\vec x| + |\vec y|$.
    \end{proof}
    
    \subsection{Orthonormal bases}

    We will consider specifically $\RR^3$ in this section. Let $\vec e_1$, $\vec e_2$, $\vec e_3$ be orthonormal, that is they are unit vectors and mutually perpendicular.

    This is an example of an \emph{orthonormal basis} -- any element of $\RR^3$ can be written as a linear combination of the elements of the basis

    \begin{definition}
        A set of vectors $\{\vec e_1, \vec e_2, \dots, \vec e_n\}$ is \emph{orthonormal} if
        \begin{itemize}
            \item They are all unit vectors -- $|\vec e_i| = 1$ for any $1 \leq i \leq n$;
            \item They are mutually perpendicular --
            
            
        \end{itemize}
        
        
    \end{definition}
    
    \begin{definition}
        A \emph{basis} for a vector space $V$ is a set of vectors such that any vector in $V$ can be uniquely represented as a linear combination of these vectors.
    \end{definition}

    \begin{definition}
        An \emph{orthonormal basis} is a basis all of whose vectors are orthonormal.
    \end{definition}
    
    Now let's say we have an orthonormal basis $\vec e_1$, $\vec e_2$, $\vec e_3$ for $\RR^3$, so we can represent every $\vec a \in \RR^3$ uniquely as a linear combination $a_1 \vec e_1 + a_2 \vec e_2 + a_3 \vec e_3$. It now becomes possible to use normal vector notation -- we write
    \begin{align*}
        \vec a &= (a_1, a_2, a_3) && \text{(row vector)} \\
        &= \begin{pmatrix}
         a_1 \\
         a_2 \\
         a_3 \\
        \end{pmatrix} && \text{(column vector)}
    \end{align*}
    
    Now for $\vec a$, $\vec b \in \RR^3$,
    \begin{align*}
        \vec a \cdot \vec b &= (a_1\vec e_1 + a_2 \vec e_2 + a_3\vec e_3) \cdot (b_1\vec e_1 + b_2 \vec e_2 + b_3\vec e_3) \\
        &= a_1b_1 + a_2b_2 + a_3b_3 && \text{by orthonormal basis}
    \end{align*}
    
    In particular, $\vec a \cdot \vec a = a_1^2 + a_2^2 + a_3^2$ -- the \emph{Pythagoras rule}.

    The \emph{canonical basis} of $\RR^3$ is such that
    \[
        \vec e_1 = (1, 0, 0) \qquad \vec e_2 = (0, 1, 0) \qquad \vec e_3 = (0, 0, 1)
    \]
    Then we write $\vec e_1 = \ihat$, $\vec e_2 = \jhat$, $\vec e_3 = \khat$.

    \subsection{Vector/cross product}

    Note: the vector product is defined only on $\RR^3$

    \begin{definition}[Vector product]
        Consider $\vec a$, $\vec b \in \RR^3$. Their vector product is given by
        \[
            \vec a \times \vec b = |\vec a||\vec b| \hvec n \sin \theta 
        \]
        where $\hvec n$ is a unit vector that is perpendicular to both $\vec a$ and $\vec b$, in a direction according to the right hand rule. We can also use the notation $\vec a \wedge \vec b$.
    \end{definition}

    \begin{remark}
        \begin{itemize}
            \item $\hvec n$ is not defined if $\vec a \parallel \vec b$, but then we say $\vec a \times \vec b = \vec 0$ as suggested by the $\sin \theta$ term;
            \item $\sin \theta$ is not defined if $|\vec a| = 0$ or $|\vec b| = 0$, but then we say $\vec a \times \vec b = \vec 0$.
        \end{itemize}
        
        
    \end{remark}
    
    

    \begin{remark}
        Note if we were to measure $\theta$ in the other direction, we would obtain $-\hvec n$ instead. From this it follows that $\vec a \times \vec b = - \vec b \times \vec a$.
    \end{remark}


    
    
    
    DIAGRAM: $\vec a$, $\vec b$ in a plane, $\vec b$ `above' $\vec a$, angle $\theta$ between them, two options for $\hvec n$, directed angle from $\vec a$ to $\vec b$, we use $\hvec n$ above the plane

    \begin{itemize}
        \item \emph{Anticommutativity.} $\vec a \times \vec b = -\vec b \times \vec a$.
        \item \emph{Alternating property} $\vec a \times \vec a = 0$.
        \item $\vec a \times \vec b = 0$ if and only if $\vec a \parallel \vec b$.
        \item $(\lambda \vec a) \times \vec b = \lambda (\vec a \times \vec b) = \vec a \times (\lambda \vec b)$ for a real parameter $\lambda$
        \item $\vec a \times (\vec b + \vec c) = \vec a \times \vec b + \vec a \times \vec c$
        \item $\vec a \cdot (\vec a \times \vec b) = \vec b \cdot (\vec a \times \vec b) = 0$.
    \end{itemize}
    
    Here are some interpretations of the cross product:

    \begin{itemize}
        \item $\vec a \times \vec b$ is the `vector area' of the parallelogram spanned by $\vec a$ and $\vec b$, and half of this is the area of the triangle.
        
    
        DIAGRAM: vectors a, b, b above a, angle theta, perpendicular dropped from b to a, dotted lines extend to parallelogram, norm of a cross b

        DIAGRAM: O, A, B, a pos vector of A, b pos vector of B, area of triangle.

    \end{itemize}
    
    \lecture{Lecture 5 -- 18/10/25}

    \begin{itemize}[resume]
        \item Fix $\vec a$ and consider some $\vec x$ with $\vec x \perp \vec a$. Then $\vec a \times \vec x$ gives $\vec x$ rotated by $\pi/2$ in a plane orthogonal to $\vec a$, and then scaled by $|\vec a|$.
        
        DIAGRAM: plane containing $\vec x$, $\vec a \times \vec x$, $\vec a$ pointing up out of the plane, three right angles, bottom one labelled $\pi/2$.
    \end{itemize}
    
    Observe the following identities:
    \begin{align*}
        \ihat \times \jhat &= \khat = -\jhat \times \ihat \\
        \jhat \times \khat &= \ihat = -\khat \times \ihat \\
        \khat \times \ihat &= \jhat = -\ihat \times \khat
    \end{align*}
    
    Now, for vectors $\vec a = (a_1, a_2, a_3)$, $\vec b = (b_1, b_2, b_3)$, we can express the cross product as follows:
    \begin{align*}
        \vec a \times \vec b =
        \begin{pmatrix}
         a_2b_3-a_3b_2 \\
         a_3b_1-a_1b_3 \\
         a_1b_2-a_2b_1 \\
        \end{pmatrix}
    \end{align*}
    
    This can also be viewed as the `determinant' of the following `matrix':
    \begin{align*}
        \vec a \times \vec b = \begin{vmatrix}
         \ihat & \jhat & \khat \\
         a_1 & a_2 & a_3 \\
         b_1 & b_2 & b_3 \\
        \end{vmatrix}
    \end{align*}
    
    \subsection{Triple products}

    \subsubsection{Scalar triple product}

    \begin{definition}
        
        For three vectors $\vec a$, $\vec b$, $\vec c \in \RR^3$, we define the \emph{scalar triple product}
        \[
            [\vec a, \vec b, \vec c] = \vec a \cdot (\vec b \times \vec c)
        \]

    \end{definition}
    
    The following identities hold:
    \begin{align*}
        \vec a \cdot (\vec b \times \vec c) = \vec b \cdot (\vec c \times \vec a) = \vec c \cdot (\vec a \times \vec b) \\
        = -\vec a \cdot (\vec c \times \vec b) = -\vec c \cdot (\vec b \times \vec a) = -\vec b \cdot (\vec a \times \vec c)
    \end{align*}
    
    DIAGRAM: vectors $\vec a$, $\vec b$, $\vec c$, $\vec a \times \vec b$ labelled, angle $\theta$ between $\vec a \times \vec b$ and $\vec c$, parallelepiped filled in with dotted lines. LABEL: $|\vec c \cdot (\vec a \times \vec b)|$ is the volume of the parallelepiped, given by the area of the parallelogram base times the perpendicular height.

    Importantly, $\vec a \cdot (\vec b \times \vec c)$ is a \emph{signed} volume -- that is,
    \begin{itemize}
        \item If $\vec a \cdot (\vec b \times \vec c) > 0$, then $\vec a$, $\vec b$, $\vec c$ constitute a \emph{right-handed set}.
        \item $\vec a \cdot (\vec b \times \vec c)$ if and only if $\vec a$, $\vec b$, $\vec c$ are \emph{coplanar}.
    \end{itemize}
    
    \subsubsection{Vector triple product}

    \begin{definition}
        For $\vec a$, $\vec b$, $\vec c \in \RR^3$, we define the \emph{vector triple product}
        \[  
            \vec a \times (\vec b \times \vec c) = (\vec a \cdot \vec c) \vec b - (\vec a \cdot \vec b) \vec c
        \]
        It is not associative!
    \end{definition}
    
    
    \subsection{Lines, planes and vector equations}

    \subsubsection{Lines}

    Any point $\vec r$ on a line through $\vec a$ with direction $\vec u$ has position vector

    \[  
        \vec r = \vec a + \lambda \vec u 
    \]
    for $\lambda \in \RR$ -- \emph{parametric form}.

    We can also express a line in the form
    \begin{align*}
        \vec u \times (\vec r - \vec a) &= \vec 0 \\
        \iff \vec u \times \vec r &= \vec b \text{ constant } (\vec b = \vec u \times \vec a).
    \end{align*}

    
    DIAGRAM: origin, $\vec a$, $\vec u$ labelled starting at $\vec a$ along with $\lambda \vec u$, $\vec r$ labelled going to the tip

    \subsubsection{Planes}

    Any point on a plane through $\vec a$ with directions $\vec u$ and $\vec v$ (non-parallel) has position vector
    \[
        \vec r = \vec a + \lambda \vec u + \mu \vec v
    \]
    for $\lambda$, $\mu \in \RR$ -- parametric form.

    If we take $\vec n = \vec u \times \vec v$, we have the plane's \emph{normal vector}, and so we can obtain the non-parametric form

    \[  
        (\vec r - \vec a) \cdot \vec n = 0 \implies \vec r \cdot \vec n = \vec r \cdot \vec a = \text{const}
    \]

    DIAGRAM: plane with all of these things labelled.

    The component of $\vec r$ along $\vec n$ is

    \[
        \hvec n \cdot \vec r = \frac{\vec n \cdot \vec r}{|\vec n|} = \frac{\vec n \cdot \vec a}{|\vec n|}
    \]

    and this quantity, independent from $\vec r$, is the perpendicular distance from the origin to the plane.

    \begin{proposition}[Intersection of line and plane]
        Given a line $\vec u \times \vec r = \vec u \times \vec a$ and a plane $\vec n \cdot \vec r = \vec n \cdot \vec b$ satisfying $\vec u \cdot \vec n \neq 0$, their intersection is given by
        \[
            \vec r = \frac{(\vec b \cdot \vec n)\vec u - (\vec a \times \vec u) \times \vec n}{\vec u \cdot \vec n}
        \]
    \end{proposition}
    
    \begin{proof}
        Suppose $\vec r$ lies on the line, so
        \begin{align*}
            \vec a \times \vec u &= \vec r \times \vec u \\
            \implies (\vec a \times \vec u) \times \vec n &= (\vec r \times \vec u) \times \vec n \\
            &= (\vec r \cdot \vec n) \vec u - (\vec u \cdot \vec n) \vec r && \text{(vector triple product)} \\
            &= (\vec b \cdot \vec n) \vec u - (\vec u \cdot \vec n) \vec r && \text{(equation of plane)} \\
            \implies (\vec u \cdot \vec n) \vec r &= (\vec b \cdot \vec n) \vec u - (\vec a \times \vec u) \times \vec n
        \end{align*}
        Hence, if $\vec u \cdot \vec n \neq 0$,
        \[
            \vec r = \frac{(\vec b \cdot \vec n)\vec u - (\vec a \times \vec u) \times \vec n}{\vec u \cdot \vec n}
        \]
        as required. Otherwise, the line and plane are either parallel and never intersect, or the line lies in the plane.
        
    \end{proof}
    
    
    

    \lecture{Lecture 6 -- 21/10/25}

    \begin{proposition}[Shortest distance between two lines]
        Consider two non-intersecting lines:
        \begin{align*}
            \ell_1 &: \vec u_1 \times (\vec r - \vec a_1) = \vec 0 \\
            \ell_2 &: \vec u_2 \times (\vec r - \vec a_2) = \vec 0
        \end{align*}
        The shortest distance between them is given by
        \begin{align*}
            s = \left|(\vec a_1 - \vec a_2) \cdot \frac{\vec u_1 \times \vec u_2}{|\vec u_1 \times \vec u_2|}\right|
        \end{align*}
    \end{proposition}
    
    \begin{proof}[Sketch proof]
        DIAGRAM: both lines labelled, starting from origin, etc.

        The shortest distance between $\ell_1$ and $\ell_2$ is attained along the line perpendicular to both of them, and hence parallel to the cross product of the direction vectors.

        Using this insight, we compute the distance by projecting the difference $\vec a_1 - \vec a_2$ onto $\vec u_1 \times \vec u_2$.
    \end{proof}
    
    \subsubsection{Vector equations}

    Goal: solve
    \begin{align*}
        \vec r + \vec a \times (\vec b \times \vec r) = \vec c && (*)
    \end{align*}
    for some specified vectors $\vec a$, $\vec b$, $\vec c$.

    First, we rewrite using the vector triple product identity:
    \begin{align*}
        \vec r + (\vec a \cdot \vec r)\vec b - (\vec a \cdot \vec b) \vec r &= \vec c \\
        \implies \vec a \cdot \left(\vec r + (\vec a \cdot \vec r)\vec b - (\vec a \cdot \vec b) \vec r\right) &= \vec a \cdot \vec c \\
        \implies \vec a \cdot \vec r + (\vec a \cdot \vec r)(\vec a \cdot \vec b) - (\vec a \cdot \vec b)(\vec a \cdot \vec r) &= \vec a \cdot \vec c \\
        \implies \vec a \cdot \vec r &= \vec a \cdot \vec c.
    \end{align*}
    Thus
    \begin{align*}
        \vec r + (\vec a \cdot \vec c)\vec b - (\vec a \cdot \vec b) \vec r &= \vec c \\
        \implies \vec r (1 - \vec a \cdot \vec b) &= \vec c - (\vec a \cdot \vec c) \vec b
    \end{align*}
    So if $\vec a \cdot \vec b \neq 1$, we can divide by $1 - \vec a \cdot \vec b$, to get a unique solution
    \[
        \vec r = \frac{\vec c - (\vec a \cdot \vec c) \vec b}{1 - \vec a \cdot \vec b}
    \]
    in other words, a point.

    Otherwise, if $\vec a \cdot \vec b = 1$ and $\vec c - (\vec a \cdot \vec c)\vec b \neq 0$, there are clearly no solutions. In the case where $\vec c - (\vec a \cdot \vec c)\vec b = 0$, the solutions are given by any $\vec r$ satisfying
    \[
        \vec a \cdot \vec r = \vec a \cdot \vec c
    \]
    in other words, a plane. 

    In fact it is necessary to verify that these conditions are sufficient as well as necessary.

    \subsubsection{Index notation and summation conventions}

    Consider an orthonormal right-handed basis $\vec e_1$, $\vec e_2$, $\vec e_3$ (not necessarily the canonical basis).

    \begin{definition}[Kronecker delta]
        The \emph{Kronecker delta} is given by
        \begin{align*}
            \delta_{ij} = 
            \begin{cases}
                1 & i = j \\
                0 & i \neq j
            \end{cases}
        \end{align*}
        Clearly it is symmetric.
    \end{definition}
    Now, we are able to represent our orthonormality condition: $\vec e_i \cdot \vec e_j = \delta_{ij}$. This allows a simple derivation of the coordinate dot product:
    \[
        \vec a \cdot \vec b = \sum_{i,j=1}^3 \delta_{ij}a_ib_j = \sum_{i=1}^3 a_ib_i
    \]
    \begin{definition}[Levi-Civita epsilon]
        The \emph{Levi-Civita epsilon} is given by
        \begin{align*}
            \lv{ijk} =
            \begin{cases}
                1 & (i, j, k) \text{ is an even permutation of } (1, 2, 3) \\
                -1 & (i, j, k) \text{ is an odd permutation of } (1, 2, 3) \\
                0 & \text{o/w}
            \end{cases}
        \end{align*}
        Simply, $\lv{123} = \lv{231} = \lv{312} = 1$, $\lv{132} = \lv{213} = \lv{321} = -1$, and 0 otherwise.

        Now it is \emph{antisymmetric}.
    \end{definition}
    Now we can also derive the coordinate cross product:
    \[
        \vec a \times \vec b = \sum_{i,j,k=1}^3 \lv{ijk} a_i b_j \vec e_k
    \]

    \begin{definition}[Einstein's summation convention]
        Indices that appear twice in an expression are assumed to be summed over -- to simplify notation, we omit the sigma notation.
    \end{definition}
    
    \begin{example}
        \begin{enumerate}[label=\protect\circled{\arabic*}]
            \item $a_i \delta_{ij} = a_j$ means $\sum_i a_i \delta_{ij} = a_j$.
            \item $\vec a \cdot \vec b = \delta_{ij} a_i b_j = a_ib_i$ means $\vec a \cdot \vec b = \sum_{i,j} \delta_{ij} a_i b_j = \sum_i a_i b_i$.
            \item $(\vec a \times \vec b)_i = \lv{ijk} a_j b_k$
            \item $\vec a \cdot (\vec b \times \vec c) = \lv{ijk} a_i b_j c_k$
            \item $\delta_{ii} = 3$.
        \end{enumerate}
        
        
    \end{example}
    
    \begin{proposition}[Identities of Kronecker delta + Levi-Civita epsilon]
        \begin{itemize}
            \item $\lv{ijk} \lv{pqr} = \delta_{ip} \delta_{jq} \delta_{kr} - \delta_{jp} \delta_{iq} \delta_{kr} + \delta_{jp} \delta_{kq} \delta_{ir} - \delta_{kp} \delta_{jq} \delta_{ir} + \delta_{kp} \delta_{iq} \delta_{ji} - \delta_{ip} \delta_{kq} \delta_{ji}$;
            \item $\lv{ijk}\lv{pqk} = \delta_{ip} \delta_{jq} - \delta_{iq} \delta_{jp}$;
            \item $\lv{ijk} \lv{pjk} = 2 \delta_{ip}$;
            \item $\lv{ijk} \lv{ijk} = 6$.
        \end{itemize}
    \end{proposition}

    These identities are powerful for proving statements involving vectors very quickly.
    
    \begin{proposition}[Vector triple product]
        $\vec a \times (\vec b \times \vec c) = (\vec a \cdot \vec c) \vec b - (\vec a \cdot \vec b) \vec c$
    \end{proposition}
    
    \begin{proof}
        \begin{align*}
            [\vec a \times (\vec b \times \vec c)]_i &= \lv{ijk} a_j (\vec b \times \vec c)_k \\
            &= \lv{ijk}a_j \lv{kpq} b_p c_q \\
            &= \lv{ijk} \lv{pqk} a_j b_p c_q \\
            &= (\delta_{ip} \delta_{jq} - \delta_{iq} \delta_{jp}) a_j b_p c_q \\
            &= a_j b_i c_j - a_j b_j c_i \\
            &= (\vec a \cdot \vec c) b_i - (\vec a \cdot \vec b) c_i \\
        \end{align*}

    \end{proof}
    
    \lecture{Lecture 7 -- 23/10/25}

    \subsection{Spherical trigonometry}

    First, let's have a proposition.

    \begin{proposition}
        For vectors $\vec a$, $\vec b$, $\vec c \in \RR^3$,
        \[
            (\vec a \times \vec b) \cdot (\vec b \times \vec c) = (\vec a \cdot \vec b)(\vec b \cdot \vec c) - (\vec a \cdot \vec c)|\vec b|^2
        \]
        \label{dot-of-crosses-identity}
    \end{proposition}
    
    \begin{proof}
        \begin{align*}
            \mathrm{LHS} &= (\vec a \times \vec b)_i \cdot (\vec b \times \vec c)_i && \text{Einstein: summing over } i\\
            &= \lv{ijk}a_jb_k \cdot \lv{ipq}b_pc_q \\
            &= \lv{ijk}\lv{ipq}a_jb_kb_pc_q \\
            &= \lv{jki}\lv{pqi}a_jb_kb_pc_q \\
            &= (\delta_{jp}\delta_{kq} - \delta_{jq}\delta_{kp})a_jb_kb_pc_q \\
            &= a_jb_kb_jc_k - a_jb_kb_kc_j \\
            &= \underbrace{a_jb_j} \underbrace{b_kc_k} - \underbrace{a_jc_j} \underbrace{b_kb_k} \\
            &= (\vec a \cdot \vec b)(\vec b \cdot \vec c) - (\vec a \cdot \vec c) |\vec b|^2
        \end{align*}        
    \end{proof}
    
    Now consider a unit sphere with centre $O$, and points $A$, $B$, $C$ on the surface of the sphere with position vectors $\vec a$, $\vec b$, $\vec c$.

    DIAGRAM: sphere, spherical triangle vertices A, B, C, angle alpha at BAC, length A to B labelled $\delta(A, B)$

    The distance from $A$ to $B$, $\delta(A, B)$ is the arc length of the sphere and so, since we are on a unit sphere, it is also the angle in radians between $A$ and $B$:

    \[
        \angle AOB = \delta(A, B)
    \]

    It follows that
    \begin{align*}
        \vec a \cdot \vec b &= \cos \angle AOB && \text{since we're on a unit sphere} \\
        &= \cos \delta(A, B)
    \end{align*}
    and likewise,
    \[
        |\vec a \times \vec b| = \sin \delta(A, B)
    \]

    Now we have what we need for a new proposition.

    \begin{proposition}[Spherical law of cosines]
        For a unit sphere, we have
        \[
            \cos \delta(B, C) = \cos \delta(A, B) \cos \delta(A, C) + \sin \delta(A, B) \sin \delta(A, C) \cos \alpha
        \]
        where $\alpha$ is the angle between $AB$ and $AC$.
    \end{proposition}
    
    \begin{proof}
        First, observe that $\alpha$, the angle between $AB$ and $AC$, can be viewed as the angle between the planes through $OAB$ and $OAC$. Hence, it is equal to the angle between their normal vectors, and so
        \begin{align*}
            \cos \alpha &= \frac{(\vec a \times \vec b) \cdot (\vec a \times \vec c)}{|\vec a \times \vec b| \cdot |\vec a \times \vec c|} \\
            &= -\frac{(\vec b \times \vec a) \cdot (\vec a \times \vec c)}{|\vec a \times \vec b| \cdot |\vec a \times \vec c|} \\
            &= \frac{(\vec b \cdot \vec c)|\vec a|^2 - (\vec b \cdot \vec a)(\vec a \cdot \vec c)}{|\vec a \times \vec b| \cdot |\vec a \times \vec c|} && \text{by \cref{dot-of-crosses-identity}} \\
            &= \frac{\cos \delta(B, C) - \cos \delta(B, A) \cos \delta(A, C)}{\sin \delta(A, B) \sin \delta(A, C)}
        \end{align*}
        and the result is then a simple rearrangement.
    \end{proof}
    
    But how can we actually express a sphere in $\RR^3$? A sphere in $\RR^3$ with centre $O$ and radius $r \in \RR^+$ is given by
    
    \[
        \Sigma = \{\vec x \in \RR^3: |\vec x| = r\}
    \]
    
    More generally, a hypersphere in $\RR^n$ with centre $\vec a$ and radius $r$ as before is given by

    \[
        \Sigma = \{\vec x \in \RR^n: |\vec x - \vec a| = r\}
    \]
    
    DIAGRAM: 0 drawn, $\vec a$ drawn, sphere drawn, point on sphere $\vec x$.

    \subsection{$\RR^n$}

    \begin{definition}
        $\RR^n$ is the vector space defined as follows:
        \begin{itemize}
            \item $\RR^n = \{\vec x = (x_1, \dots, x_n): x_i \in \RR\}$.
            \item \emph{Addition}. Given $\vec x = (x_1, \dots, x_n)$ and $\vec y = (y_1, \dots, y_n)$, $\vec x + \vec y$ is given by $(x_1 + y_1, \dots, x_n + y_n)$.
            \item \emph{Scalar multiplication}. Given $\vec x = (x_1, \dots, x_n)$ and $\lambda \in \RR$, $\lambda \vec x = (\lambda x_1, \dots, \lambda x_n)$.
        \end{itemize}
        One can verify that these satisfy the relevant axioms.

        We also introduce some other important notions for $\RR^n$:
        \begin{itemize}
            \item \emph{Canonical basis}. The canonical basis $\{\vec e_1, \dots, \vec e_n\}$ is given by:
            \begin{align*}
                \vec e_1 &= (1, \dots, 0) \\
                \vdots \\
                \vec e_n &= (0, \dots, 1)
            \end{align*}
            \item \emph{Inner product}. The inner product of $\vec x$ and $\vec y$ is given by
            \[
                \vec x \cdot \vec y = \sum_i x_iy_i
            \]
            It is easy to check that
            \[
                \vec e_i \cdot \vec e_j = \delta_{ij}
            \]
            and so a given component of $\vec x = (x_1, \dots, x_n)$ is given by
            \[
                x_i = \vec e_i \cdot \vec x
            \]
            If we instead write our vectors as column vectors, then for $\vec x$, $\vec y \in \RR^n$, $\vec x^\intercal$ and $\vec y^\intercal$ are column vectors, and so, using the matrix multiplication which we will define later,
            \[
                \vec x \cdot \vec y = \vec x^\intercal \vec y
            \]
        \end{itemize}
    \end{definition}
    
    \subsubsection{Generalised index notation and summation convention}

    How do our notations and conventions generalise to $\RR^n$? The Kronecker delta is fairly obvious and in fact we already used it to represent the dot product of our basis vectors. The Levi-Civita epsilon is similar:

    \begin{definition}[Generalised Levi-Civita epsilon]
        The generalised Levi-Civita epsilon is given by
        \begin{align*}
            \lv{a_1a_2\dots a_n} =
            \begin{cases}
                1 & (a_1, a_2, \dots, a_n) \text{ is an even permutation of } (1, 2, \dots, n) \\
                -1 & (a_1, a_2, \dots, a_n) \text{ is an odd permutation of } (1, 2, \dots, n) \\
                0 & \text{o/w}
            \end{cases}
        \end{align*}
        We again have antisymmetry -- swapping any two indices negates the sign.
    \end{definition}

    In $\RR^2$, this gives us a new scalar product:

    \[
        [\vec a, \vec b] = \lv{ij}a_ib_j = a_1b_2 - a_2b_1
    \]

    Geometrically, it represents the signed area of a parallelogram spanned by $\vec a$ and $\vec b$.

    DIAGRAM: $\vec a$, $\vec b$ from the same point, dotted lines complete a parallelogram.

    This is essentially analogous to $[\vec a, \vec b, \vec c] = \vec a \cdot (\vec b \times \vec c)$.

    \subsection{$\CC^n$}

    \begin{definition}
        $\CC^n$ is in large part analogous to $\RR^n$.
        \begin{itemize}
            \item $\CC^n = {\vec z = (z_1, \dots, z_n): z_i \in \CC}$
            \item \emph{Addition}. Component-by-component, as for $\RR^n$.
            \item \emph{Scalar multiplication}. Again, component-by-component, as for $\RR^n$.
        \end{itemize}
        Onto the basis and inner product, which are of note.
        \begin{itemize}
            \item \emph{Canonical basis}. This depends on whether we allow our scalars to be real or complex.
            \begin{itemize}
                \item \emph{Real scalars only}. Then we have a canonical basis $\{\vec e_1, \dots, \vec e_n, \vec f_1, \dots, \vec f_n\}$ of size $2n$, given by
                \begin{align*}
                    \vec e_1 &= (1, \dots, 0) && \vec f_1 = (i, \dots, 0) \\
                    \vdots \\
                    \vec e_n &= (0, \dots, 1) && \vec f_n = (0, \dots, i)
                \end{align*}
                Here, $\CC^n$ is a real vector space (even though its vectors have complex number components), and is the same as $\RR^{2n}$.
                \item \emph{Complex scalars}. Now we can introduce our complex numbers by scalar multiplication, so we don't need separate basis vectors for $i$. Thus our basis for $\RR^n$, $\{\vec e_1, \dots, \vec e_n\}$ is also a basis for $\CC^n$ as a \emph{complex vector space} with dimension $n$.
            \end{itemize}
            \item \emph{Inner product}. The inner product of $\vec z$ and $\vec w$ is given by
            \[
                (\vec z, \vec w)= \sum_j \overline{z_j} w_j
            \]
        \end{itemize}
    \lecture{Lecture 8 -- 25/10/25}
            \begin{itemize}[resume]       
        
            \item Properties of the inner product:
            \begin{itemize}
                \item \emph{Hermitianity}. $(\vec w, \vec z) = \overline{(\vec z, \vec w)}$ \\
                \item \emph{Linearity/antilinearity}.
                \begin{align*}
                    (\vec z, \lambda \vec w + \lambda ' \vec w') &= \lambda (\vec z, \vec w) + \lambda ' (\vec z, \vec w') \\
                    (\mu \vec z + \mu' \vec z', \vec w) &= \bar \mu (\vec z, \vec w) + \overline{\mu'}(\vec z', \vec w)
                \end{align*}
                \item \emph{Positive-definiteness}.
                \[
                    (\vec z, \vec z) = \sum_j |\vec z_j|^2 \geq 0
                \]
                with equality if and only if $\vec z = 0$.
                \item \emph{Norm}. $|\vec z| = \sqrt{(\vec z, \vec z)}$.
            \end{itemize}
        \end{itemize}
    \end{definition}

    \begin{definition}[Orthogonality in $\CC^n$]
        We say $\vec z$, $\vec w \in \CC^n$ are \emph{orthogonal} if
        \[
            (\vec z, \vec w) = 0
        \]
    \end{definition}
    
    \begin{remark}
        It now becomes clear that our standard basis for $\CC^n$ is in fact an orthonormal basis.
    \end{remark}
    
    \subsubsection{From complex to real inner products}

    In the case $n=1$, $z$, $w \in \CC$,
    \[  
        (z, w) = \bar z w
    \]
    Then, put $z = a_1 + ia_2$, $w = b_1 + ib_2$, and identify these with $\vec a = (a_1, a_2)$, $\vec b = (b_1, b_2) \in \RR^2$. Then
    \[  
        (z, w) = \vec a \cdot \vec b + i [\vec a, \vec b]
    \]
    So one can recover from the inner product in $\CC$ both of our scalar products in $\RR^2$.

    \subsection{General vector spaces}

    \begin{definition}[General vector space]
        We can define a vector space over any field, not just $\RR$ or $\CC$.
        \begin{itemize}
            \item If our scalars are from $\RR$ or $\CC$, then we have a \emph{real vector space} or \emph{complex vector space} respectively.
        \end{itemize}
    \end{definition}

    \begin{definition}
        Given a real vector space $V$, $\vec v_1, \dots, \vec v_r \in V$, a \emph{linear combination} of them is given by
        \[  
            \lambda_1 \vec v_1 + \dots + \lambda r \vec v_r \in V
        \]
        where the $\lambda_i \in \RR$.
    \end{definition}
    
    \begin{definition}
        The \emph{span} of $\vec v_1, \dots, \vec v_r \in V$, is the set of all linear combinations of $\vec v_1, \dots, \vec v_r$ -- that is,
        \[
            \Span\{\vec v_1, \dots, \vec v_r\} = \left\{\sum_i \lambda_i \vec v_i: \lambda_i \in \RR\right\}
        \]
    \end{definition}
    
    \begin{definition}
        A \emph{subspace} of $V$ is a subset of $V$ that is itself a vector space. Equivalently, a non-empty subset $U \subseteq V$ is a subspace if, for any $\vec v$, $\vec w \in U$, and any $\lambda$, $\mu \in \RR$, $\lambda \vec v + \mu \vec w \in U$.
    \end{definition}
    
    \begin{proposition}
        The span of any set of vectors is a subspace.
    \end{proposition}
    
    \begin{remark}
        The two trivial subspaces of any vector space $V$ are $V$ and $\{0\}$.
    \end{remark}
    
    \subsubsection{Linear independence and dependence}
    
    In a general vector space, consider the following equation for some fixed $\vec v_1, \dots, \vec v_n$:
    
    \[
        \lambda_1 \vec v_1 + \dots + \lambda_n \vec v_n = \vec 0
    \]
    
    where the $\lambda_i$ come from the field in question.

    \begin{definition}
        If the only solution to this equation is $\lambda_1 = \dots = \lambda_n = 0$, the vectors $\vec v_1, \dots, \vec v_n$ are \emph{linearly independent}, otherwise \emph{linearly dependent}.
    \end{definition}
    
    \begin{example}
        \begin{itemize}
            \item Consider $\vec v_1 = (1, 0)$, $\vec v_2 = (0, 1)$, $\vec v_3 = (0, 2) \in \RR^2$. Clearly this set is linearly dependent -- $0 \cdot \vec v_1 + 2 \cdot \vec v_2 - 1 \cdot \vec v_3 = 0$. However, $\{\vec v_1, \vec v_2\}$ do comprise a linearly independent set.
            \item Any set containing $\vec 0$ is linearly dependent.
            \item In $\RR^3$, $\vec a$, $\vec b$, $\vec c$ are linearly independent if and only if $\vec a \cdot (\vec b \times \vec c) \neq 0$.
        \end{itemize}
    \end{example}
    
    \subsubsection{General inner products}

    The general inner product satisfies the same properties as the inner product in $\CC^n$. For completeness, let us restate them.

    \begin{itemize}
        \item $(\vec v, \vec w) \in \RR$ or $\CC$ as relevant;
        \item \emph{Hermitianity}. $(\vec v, \vec w) = \overline{(\vec v, \vec w)}$. Observe that, in a real vector space, hermitianity reduces to symmetry.
        \item \emph{Linearity/antilinearity}.
        \begin{align*}
            (\vec z, \lambda \vec w + \lambda ' \vec w') &= \lambda (\vec z, \vec w) + \lambda ' (\vec z, \vec w') \\
            (\mu \vec z + \mu' \vec z', \vec w) &= \bar \mu (\vec z, \vec w) + \overline{\mu'}(\vec z', \vec w)
        \end{align*}
        and again it is clear that we get linearity in the real case.
        \item \emph{Positive-definiteness}.
        \[
            (\vec v, \vec v) \geq 0
        \]
        with equality if and only if $\vec v = \vec 0$.
    \end{itemize}
    
    All of our previously seen inner products satisfy these axioms, which are powerful in that they will allow us to prove general statements about loads of different vector spaces, without explicitly defining what the inner product has to be.

    \begin{definition}
        As before, we say two vectors $\vec v$, $\vec w$ are \emph{orthogonal} if $\vec v \cdot \vec w = 0$.
    \end{definition}
    
    \begin{proposition}
        If $\vec v_1, \dots, \vec v_n$ are non-zero and orthogonal, then they are linearly independent.
    \end{proposition}
    
    \begin{proof}
        Suppose $\sum_{i} \lambda_i \vec v_i = 0$. Then select some $j$:
        \begin{align*}
            0 &= \left(\vec v_i, \sum_i \lambda_i \vec v_i\right) \\
            &= \sum_i \lambda_i (\vec v_j, \vec v_i) && \text{by linearity} \\
            &= \lambda_j (\vec v_j, \vec v_j) && \text{by orthogonality} \\
            &= \lambda_j |\vec v_j|^2            
        \end{align*}
        By positive-definiteness and the assumption that $\vec v_i \neq \vec 0$, $|\vec v_j|^2 \neq 0$, so $\lambda_j = 0$. Performing this argument for all $j$ completes the proof.
    \end{proof}
    
    \subsubsection{Bases and dimensions}
    
    \begin{definition}
        A basis $\mathcal B = \{\vec e_1, \dots, \vec e_n\}$ of a vector space $V$ is such that:
        \begin{itemize}
            \item $\mathcal B$ is linearly independent;
            \item $\Span \mathcal B = V$.
        \end{itemize}
        Importantly, each vector in $V$ can only be expressed as a linear combination of basis elements in one way -- otherwise, subtracting them would contradict linear independence.

        Then, for a vector $\vec v = v_1 \vec e_1 + \dots + v_n \vec e_n$, $v_1, \dots, v_n$ are the \emph{components of $\vec v$ with respect to $\mathcal B$}.
    \end{definition}

    \begin{theorem}
        If $\mathcal B_1$ and $\mathcal B_2$ are bases for $V$, then $|\mathcal B_1| = |\mathcal B_2|$.
    \end{theorem}
    
    \begin{definition}
        The number of elements in a basis for $V$ is the \emph{dimension} of $V$.
    \end{definition}
    
    \begin{proposition}
        For $V$ a vector space of dimension $n$,
        \begin{itemize}
            \item If $Y = \{\vec w_1, \dots, \vec w_m\}$ spans $V$, and $m \geq n$, there must exist a subset of $Y$ which is a basis for $V$.
            \item If $X = \{\vec u_1, \dots, \vec u_k\}$ is linearly independent, and $k \leq n$, there must exist a superset of $X$ which is a basis for $V$.
        \end{itemize}
    \end{proposition}

    \lecture{Lecture 9 -- 28/10/25}
    
    \section{Matrices}

    \subsection{Linear maps}

    \begin{definition} 
        Consider two vector spaces $V$ and $W$. A linear map is a function $T: V \to W$ such that $T(\lambda \vec x + \mu \vec y) = \lambda T(\vec x) + \mu T(\vec y)$ for all scalars $\lambda$, $\mu$ and vectors $\vec x$, $\vec y$.
    \end{definition}
    
    \begin{definition}
        The \emph{image} of a linear map is the same as for any function:
        \begin{itemize}
            \item if $T(\vec x) = \vec x'$, then $\vec x'$ is called the image of $\vec x$ under $T$;
            \item the image of $T$ itself, $\Im T = \{\vec x' \in W : \vec x' = T(\vec x) \text{ for some } \vec x \in V\}$.
        \end{itemize}
    \end{definition}
    
    \begin{proposition}
        $\Im T$ is a subspace of $W$.
    \end{proposition}
    
    \begin{definition}
        The \emph{kernel} of a linear map $T$ is the set $\ker T = \{\vec x \in V: T(\vec x) = \vec 0\}$ -- the set of elements of $V$ that get sent to zero in $W$.
    \end{definition}
    
    \begin{proposition}
        $\ker T$ is a subspace of $V$.
    \end{proposition}
    
    \begin{definition}[Rank and nullity]
        \begin{itemize}
            \item The dimension of the image of $T$, $\dim(\Im T)$, is the \emph{rank} of $T$.
            \item The dimension of the kernel of $T$, $\dim(\ker T)$, is the \emph{nullity} of $T$.
        \end{itemize}
    \end{definition}
    
    \begin{remark}
        Since $\ker T$ is a subspace of $V$, the nullity $\dim(\ker T) \leq \dim V$, and likewise the rank $\dim(\Im T) \leq \dim W$.
    \end{remark}
    
    \begin{example}
        Examples of linear maps include:
        \begin{enumerate}[label=\protect\circled{\arabic*}]
            \item The \emph{zero linear map} $T: V \to W$, $\vec x \mapsto \vec 0$. Clearly $\Im T = \{\vec 0\}$, $\ker T = V$.
            \item The \emph{identity linear map} $T: V \to V$, $\vec x \mapsto \vec x$. Now $\Im T = V$, $\ker T = \{\vec 0\}$.
            \item Consider $V = W = \RR^2$,
            \begin{align*}
                T: V &\to W \\
                \begin{pmatrix}
                 x_1 \\
                 x_2 \\
                \end{pmatrix} &\mapsto
                \begin{pmatrix}
                 2x_1 + x_2 \\
                 x_1 - 3x_2 \\
                \end{pmatrix}
            \end{align*}
            Can check that $\Im T = \RR^2$, $\ker T = \{\vec 0\}$.            
        \end{enumerate}
    \end{example}

    \begin{definition}[Linear combinations of linear maps]
        Let $T, S: V \to W$ be linear maps. Then a linear combination of these maps is defined as follows:
        \begin{align*}
            \alpha T + \beta S: V &\to W \\
            \vec x &\mapsto \alpha T(\vec x) + \beta S(\vec x)
        \end{align*}
    \end{definition}
    
    \begin{definition}[Compositions of linear maps]
        Let $T: V \to W$ and $S: U \to V$ be linear maps. Then the composition of these maps is defined as follows:
        \begin{align*}
            T \circ S: U &\to W \\
            \vec x &\mapsto T(S(\vec x))
        \end{align*}
    \end{definition}
    
    \begin{proposition}
        Linear combinations and compositions of linear maps are themselves linear maps.
    \end{proposition}
    
    \begin{theorem}[Rank-nullity theorem]
        For $T: V \to W$ a linear map,
        \[
            \rank(T) + \nullity(T) = \dim(V)
        \]
    \end{theorem}
    
    \begin{proof}[Proof (non-examinable!)]
        Let $n = \dim(V)$, $m = \nullity(T)$, and we know $m \leq n$. We consider cases.
        \begin{enumerate}
            \item If $m = n$, then $T$ must be the zero map, so $\Im(T) = \{\vec 0\}$ and thus $\rank(T) = 0$, which concludes the proof.
            \item If $m < n$, let $\{\vec e_1, \dots, \vec e_m\}$ be a basis for $\ker T$, then $T(\vec e_i) = 0$ for any $i$. We can extend this basis to a basis for $V$:
            \[
                \{\vec e_1, \dots, \vec e_m, \vec e_{m+1}, \dots, \vec e_n\}
            \]
            and we claim $\{T(\vec e_{m+1}), \dots, T(\vec e_n)\}$ is a basis for $\Im T$.
            \begin{enumerate}
                \item First, we show $\{T(\vec e_{m+1}), \dots, T(\vec e_n)\}$ spans $\Im T$. Take $\vec y \in \Im T$, so there exists $\vec x \in V$ with $T(\vec x) = \vec y$. Since $\vec x \in V$, we may write
                \begin{align*}
                    \vec x &= \alpha_1 \vec e_1 + \dots + \alpha_m \vec e_m + \alpha_{m+1} \vec e_{m+1 }+ \dots + \alpha_n \vec e_n \\
                    \implies \vec y = T(\vec x) &= \underbrace{\alpha_1 T(\vec e_1) + \dots + \alpha_m T(\vec e_m)}_{\vec 0} + \alpha_{m+1} T(\vec e_{m+1})+ \dots + \alpha_n T(\vec e_n) \\
                    &= \alpha_{m+1} T(\vec e_{m+1})+ \dots + \alpha_n T(\vec e_n)
                \end{align*}
                and so $\{T(\vec e_{m+1}), \dots, T(\vec e_n)\}$ spans $\Im T$ as required.
                \item It remains to show that these vectors are linearly independent. Suppose
                \begin{align*}
                    \vec 0 &= \alpha_{m+1} T(\vec e_{m+1})+ \dots + \alpha_n T(\vec e_n) \\
                    &= T(\underbrace{\alpha_{m+1} \vec e_{m+1}+ \dots + \alpha_n \vec e_n}_{\vec x})
                \end{align*}
                and so $\vec x \in \ker T$. Thus, $\vec x$ is a linear combination of $\{\vec e_1, \dots, \vec e_m\}$ -- there exist $\alpha_1, \dots, \alpha_m$ such that
                \begin{align*}
                    \alpha_1 \vec e_1 + \dots + \alpha_m \vec e_m = \vec x = \alpha_{m+1} \vec e_{m+1} + \dots + \alpha_n \vec e_n \\
                    \implies \alpha_1 \vec e_1 + \dots + \alpha_m \vec e_m - \alpha_{m+1} \vec e_{m+1} - \dots - \alpha_n \vec e_n = \vec 0
                \end{align*}
                Then, since $\{\vec e_1, \dots, \vec e_m, \vec e_{m+1}, \dots, \vec e_n\}$ form a basis of $V$, they are linearly independent, and thus $\alpha_1 = \dots = \alpha_m = \alpha_{m+1} = \dots = \alpha_n = 0$. Hence in particular $\alpha_{m+1} = \dots = \alpha_n = 0$, so $T(\vec e_{m+1}), \dots, T(\vec e_n)$ are linearly independent as required.             
            \end{enumerate}
            Hence $\Im T$ is spanned by a basis of size $n-m$, so $\rank(T) = n-m$, and hence $\rank(T) + \null(T) = n = \dim(V)$ as required.
            
        \end{enumerate}
        
        
    \end{proof}
    
    \subsection{Matrices as linear maps}

    Let $M$ be a matrix with entries $M_{ij} \in \RR$, where $i$ labels rows and $j$ labels columns. Let $T: \RR^n \to \RR^n$ such that $T(\vec x) = \vec x' = M\vec x$, where $M \vec x$ is defined by
    \begin{align*}
        x_i' = M_{ij}x_j && \text{(note summation convention on } j \text{)}
    \end{align*}
    As an example, in the $n=2$ case, we have
    \begin{align*}
        M = \begin{pmatrix}
         M_{11} & M_{12} \\
         M_{21} & M_{22} \\
        \end{pmatrix}
    \end{align*}
    and
    \begin{align*}
        M\vec x = M\begin{pmatrix}
         x_1' \\
         x_2' \\
        \end{pmatrix}
        &= \begin{pmatrix}
         M_{11} & M_{12} \\
         M_{21} & M_{22} \\
        \end{pmatrix}
        \begin{pmatrix}
         x_1 \\
         x_2 \\
        \end{pmatrix} \\
        &= \begin{pmatrix}
         M_{11}x_1 + M_{12}x_2 \\
         M_{21}x_1 + M_{22}x_2 \\
        \end{pmatrix}
    \end{align*}
    In general, we can consider the rows $\vec R_i \in \RR^n$ and $\vec C_i \in \RR^n$, and then we can write

    \[  
        M = \begin{pmatrix}
         \leftarrow & \vec R_1 & \rightarrow \\
          & \vdots &  \\
         \leftarrow & \vec R_n & \rightarrow \\
        \end{pmatrix}
        = \begin{pmatrix}
         \uparrow &  & \uparrow \\
         \vec C_1 & \cdots & \vec C_n \\
         \downarrow &  & \downarrow \\
        \end{pmatrix}
    \]
    We then have the following `identities' (by definition):
    \[
        (\vec R_i)_j = M_ij = (\vec C_j)_i
    \]
    Then, if $\{\vec e_1, \dots, \vec e_n\}$ is the canonical basis of $\RR^n$,

    \[
        T(\vec e_i) = M \vec e_i = \vec C_i
    \]

    Since $T$ is a linear map,
    \begin{align*}
        \vec x = \sum_i x_i \vec e_i &\mapsto T\left(\sum_i x_i \vec e_i\right) \\
        &= \sum_i x_i T(\vec e_i) \\
        &= x_i \vec C_i && \text{(summation convention)}
    \end{align*}
    
    From this we have the following result:

    \begin{proposition}
        The image of a linear map $T$ is the span of the columns of the matrix $M$ representing it:
        \[
            \Im T = \Im M = \Span\{\vec C_1, \dots, \vec C_n\}
        \]
    \end{proposition}
    
    Furthermore,
    \begin{align*}
        x_i' &= M_ij x_j = (\vec R_i)_j x_j = \vec R_i \cdot \vec x \\
        \implies \vec x' &= \begin{pmatrix}
         x_1' \\
         \vdots \\
         x_n' \\
        \end{pmatrix} = \begin{pmatrix}
         \vec R_1 \cdot \vec x \\
         \vdots \\
         \vec R_n \cdot \vec x \\
        \end{pmatrix}
    \end{align*}
    Thus, if $\vec x' = \vec 0$, $\vec R_i \cdot \vec x = \dots = \vec R_n \cdot \vec x = 0$. Thus we have the following result:

    \begin{proposition}
        The kernel of a linear map $T$ is the subspace of its domain which is orthogonal to all the rows of the matrix $M$ representing it:
        \[
            \ker T = \ker M = \{\vec x \in \RR^n: \vec R_i \cdot \vec x = 0 \text{ for all }i\}
        \]
    \end{proposition}
    
    \begin{example}
        Let us consider the matrices corresponding to particular linear maps.
        \begin{itemize}
            \item Taking $V = W = \RR^n$, the \emph{zero linear map} corresponds to the matrix which is all zeroes.
            \item The \emph{identity linear map} corresponds to the matrix $I$ satisfying $I_{ij} = \delta_{ij}$.
            \item Consider $V = W = \RR^3$ and the map $T: V \to W$, $\vec x \mapsto \vec x'$ given by
            \[
                \begin{cases}
                    x_1' = 3x_1 + x_2 + 5x_3 \\
                    x_2' = -x_1 - 2x_3 \\
                    x_3' = 2x_1 + x_2 + 3x_3
                \end{cases}
            \]
            Clearly the corresponding matrix is given by
            \[
                M = \begin{pmatrix}
                 3 & 1 & 5 \\
                 -1 & 0 & -2 \\
                 2 & 1 & 3 \\
                \end{pmatrix}
            \]
            and so
            \[
                \vec C_1 = \begin{pmatrix}
                 3 \\
                 -1 \\
                 2 \\
                \end{pmatrix} \quad
                \vec C_2 = \begin{pmatrix}
                 1 \\
                 0 \\
                 1 \\
                \end{pmatrix} \quad
                \vec C_3 = \begin{pmatrix}
                 5 \\
                 -2 \\
                 3 \\
                \end{pmatrix}
            \]
            and
            \begin{align*}
                \vec R_1 &= (3, 1, 5) \\
                \vec R_2 &= (-1, 0, -2) \\
                \vec R_3 &= (2, 1, 3)
            \end{align*}
            Then 
            \begin{align*}
                \Im T = \Im M &= \Span\{\vec C_1, \vec C_2, \vec C_3\} \\
                &= \Span\{\vec C_1, \vec C_2\} && (\vec C_3 = 2 \vec C_1 - \vec C_2)
            \end{align*}
            and so $\rank(T) = 2$. To find $\ker T$, we find $\vec R_2 \times \vec R_3 = (2, -1, -1)$, and note that this is also orthogonal to $\vec R_1$, meaning that $\ker T$ is the subspace this generates -- that is,
            \[
                \ker T = \ker M = \Span \{(2, -1, -1)\}
            \]
            This makes sense, since then $\nullity(T) = 1$ and the rank-nullity theorem is satisfied.
            
        \end{itemize}
        
        
    \end{example}
    
    \subsection{Geometric examples in $\RR^2$ and $\RR^3$}

    \subsubsection{$\RR^2$}

    We will consider matrices representing rotations and reflections.

    \begin{enumerate}
        \item \emph{Rotations}. Consider $\theta$ with $-\pi < \theta \leq \pi$. Then, a rotation of angle $\theta$ counter-clockwise in $\RR^2$ is represented by
        \[
            \Rot(\theta) = \begin{pmatrix}
             \cos \theta & -\sin \theta \\
             \sin \theta & \cos \theta \\
            \end{pmatrix}
        \]
        DIAGRAM: vector in $\RR^2$ before and after rotation with angle $\theta$ indicated between.
        
        We should also note that $\det(\Rot(\theta)) = 1$ for any $\theta$.
        \item \emph{Reflections}. Consider $\theta$ with $-\pi < \theta \leq \pi$. Then, a reflection of angle $\theta/2$ in $\RR^2$ is represented by
        \[
            \Rfl(\theta) = \begin{pmatrix}
             \cos \theta & \sin \theta \\
             \sin \theta & -\cos \theta \\
            \end{pmatrix}
        \]
        and now $\det(\Rfl(\theta)) = -1$ for any $\theta$.

        DIAGRAM: vector in $\RR^2$ before and after reflection in a dotted line $\theta/2$ counter-clockwise from the $x$-axis.

        Consider in particular the case $\theta = \pi/2$, then
        \[
            \Rfl(\pi/2) = \begin{pmatrix}
             0 & 1 \\
             1 & 0 \\
            \end{pmatrix}
        \]
        which has the effect of swapping the $x$ and $y$ coordinates, as we would expect.
    \end{enumerate}
    
    
    These transformations have properties including the following:

    \begin{itemize}
        \item $\Rot(\theta)\Rot(\phi) = \Rot(\theta + \phi)$
        \item $\Rfl(\theta)\Rfl(\phi) = \Rot(\theta - \phi)$
        \item $\Rot(\theta)\Rfl(\phi) = \Rfl(\phi + \theta/2)$
        \item $\Rfl(\phi)\Rot(\theta) = \Rfl(\phi - \theta/2)$
    \end{itemize}

    \begin{enumerate}[resume]
        \item \emph{Shears}. Given orthonormal vectors $\vec a$, $\vec b$, a shear $S$ with parameter $\lambda$ satisfies
        \begin{align*}
            \vec x' = S \vec x &= \vec x + \lambda \vec a (\vec x \cdot \vec b) \\
            \text{or equivalently }x_i' &= S_{ij} x_j, \; S_{ij} = \delta_{ij} + \lambda a_i b_j
        \end{align*}
        In particular,
        \begin{align*}
            S\vec a &= \vec a \\
            S\vec b &= \vec b + \lambda \vec a            
        \end{align*}
        and in fact $S \vec u = \vec u$ for any $\vec u \perp \vec b$.

        DIAGRAM: square with bottom edge vec a, left edge vec b moves to parallelogram with bottom edge vec a, left edge vec b plus lambda vec a, original vec b shown dotted.
        
    \end{enumerate}
    
    \subsubsection{$\RR^3$}
    
    \begin{enumerate}
        \item \emph{Rotations}
        \begin{enumerate}
            \item \emph{Rotations about the $z$-axis/in the $x$-$y$ plane}. Then
            \[
                \Rot(\theta) = \begin{pmatrix}
                 \cos \theta & -\sin \theta & 0 \\
                 \sin \theta & \cos \theta & 0 \\
                 0 & 0 & 1 \\
                \end{pmatrix}
            \]
            since the $z$-axis remains fixed.
            \item \emph{Rotations about an arbitrary vector}. Suppose we have a unit vector $\hvec n$ to rotate about, and our rotation takes $\vec x$ to $\vec x' = R \vec x$. Then
            \begin{align*}
                \vec x' &= (\cos \theta) \vec x + (1 - \cos \theta) (\hvec n \cdot \vec x) \hvec n + (\sin \theta) \hvec n \times \vec x \\
                R_{ij} &= (\cos \theta) \delta_{ij} + (1 - \cos \theta) n_i n_j - (\sin \theta) \lv{ijk} n_k
            \end{align*}
            To derive this, let us represent our vector as
            \[
                \vec x = \vec x_\parallel + \vec x_\perp
            \]
            where $\vec x_\parallel = (\vec x \cdot \hvec n) \hvec n$ is the component of $\vec x$ parallel to $\hvec n$ and $\vec x_\perp$ is the component of $\vec x$ perpendicular to $\hvec n$, satisfying $\vec x_\perp \cdot \hvec n = 0$.

            Then, after applying the map,
            \begin{align*}
                \vec x'_\parallel &= \vec x_\parallel \\
                \vec x'_\perp &= (\cos \theta) \vec x_\perp + (\sin \theta) \hvec n \times \vec x_\perp
            \end{align*}
            where we can see the second formula by considering the plane perpendicular to $\hvec n$ and noting $|\vec x_\perp| = |\hvec n \times \vec x|$ -- the sine and cos terms are analogous to 2D rotation. Finally, we observe $\hvec n \times \vec x_\perp = \hvec n \times \vec x$, since $\hvec n \times \vec x_\parallel = \vec 0$, and this, along with substituting $\vec x_\parallel = (\vec x \cdot \hvec n) \hvec n$ and $\vec x_\perp = \vec x - \vec x_\parallel$, recovers the formula.


            DIAGRAM: perpendicular vectors $\vec x_\perp$ and $\hvec n \times \vec x$, angle $\theta$ between $\vec x_\perp$ and $\vec x'_\perp$.

            We denote this (counter-clockwise) rotation by $\Rot(\theta, \hvec n)$.
            
            
        \end{enumerate}
        \item \emph{Reflections}. A reflection in a plane through $\vec 0$ with normal vector $\hvec n$ taking $\vec x$ to $\vec x' = H\vec x$ is given by

        \begin{align*}
            \vec x' &= \vec x - 2(\vec x \cdot \hvec n) \hvec n \\
            x_i' &=  H_{ij}x_j, \; H_{ij} = \delta_{ij} - 2n_in_j \\
            \implies H &= \begin{pmatrix}
             1-2n_1^2 & -2n_1n_2 & -2n_1n_3 \\
             -2n_1n_2 & 1-2n_2^2 & -2n_2n_3 \\
             -2n_1n_3 & -2n_2n_3 & 1-2n_3^2 \\
            \end{pmatrix}
        \end{align*}        
    \end{enumerate}
    
    \lecture{Lecture 11 -- 01/11/25}\

    \begin{enumerate}[resume]
        \item \emph{Dilations}. A dilation of factors $\alpha$, $\beta$, $\gamma$ takes
        \[
            \vec x = x_1 \vec e_1 + x_2 \vec e_2 + x_3 \vec e_3 \quad \longmapsto \quad \vec x' = \alpha x_1 \vec e_1 + \beta x_2 \vec e_2 + \gamma x_3 \vec e_3
        \]
        The matrix representing this transformation is
        \[
            M = \begin{pmatrix}
             \alpha & 0 & 0 \\
             0 & \beta & 0 \\
             0 & 0 & \gamma \\
            \end{pmatrix}
        \]
        If we want to dilate along other axes, we can compose rotations with this simplest dilation.
    \end{enumerate}
    
    \subsection{Matrices in general}

    \subsubsection{Definitions}

    Consider a linear map $T: V \to W$, $\dim V = n$, $\dim W = m$, and take two bases for $V$ and $W$ respectively:
    \begin{align*}
        \{\vec e_1, \dots, \vec e_n\} &\text{ a basis of } V \\
        \{\vec f_1, \dots, \vec f_m\} &\text{ a basis of } W
    \end{align*}
    
    We realise that $T$ is uniquely determined by where it sends the basis vectors of $V$, and we can represent these images as linear combinations of the basis vectors of $W$. So it follows that $T$ can be represented by $M$, an $m$ by $n$ array with entries $M_{ij}$, where $i = 1, \dots, m$ are the rows and $j = 1, \dots, n$ the columns.
    
    Then, any component of the image vector $x_i'$ is given by
    \[
        x_i' = \sum_j M_{ij} x_j
    \]
    and, as long as it is clear that matrix-vector multiplication is defined in this way, we can simply write
    \[
        \begin{pmatrix}
         x'_1 \\
         \vdots \\
         x'_m \\
        \end{pmatrix} =
        \begin{pmatrix}
         M_{11} & \cdots & M_{1n} \\
         \vdots &  & \vdots \\
         M_{m1} & \cdots & M_{mn} \\
        \end{pmatrix}
        \begin{pmatrix}
         x_1 \\
         \vdots \\
         x_n \\
        \end{pmatrix}
    \]
    
    If we consider another linear map $T$, represented by matrix $N$, with bases the same, then, for scalars $\alpha$ and $\beta$, the transformation
    \[
        \alpha S + \beta T
    \]
    is represented by the matrix
    \[
        \alpha M + \beta N
    \]
    with coefficients
    \[
        (\alpha M + \beta N)_{ij} = \alpha M_{ij} + \beta N_{ij}
    \]
    -- that is, scalar multiplication and addition of matrices is element-by-element.

    \begin{example}
        Take $V = M_{2 \times 2}(\RR)$ (the set of 2 by 2 real matrices) and $W = \RR^3$. It's clear that $\dim V = 4$, $\dim W = 3$. Consider the linear map
        \begin{align*}
            T: V &\to W \\
            \begin{pmatrix}
             a & b \\
             c & d \\
            \end{pmatrix} &\mapsto
            \begin{pmatrix}
             a+b \\
             c \\
             d \\
            \end{pmatrix}
        \end{align*}
        What matrix corresponds to this map? First, take a basis of $V$ and $W$:
        \begin{align*}
            V&: \left\{\vec e_1 = \begin{pmatrix}
             1 & 0 \\
             0 & 0 \\
            \end{pmatrix},\,
            \vec e_2 = \begin{pmatrix}
             0 & 1 \\
             0 & 0 \\
            \end{pmatrix},\,
            \vec e_3 = \begin{pmatrix}
             0 & 0 \\
             1 & 0 \\
            \end{pmatrix},\,
            \vec e_4 = \begin{pmatrix}
             0 & 0 \\
             0 & 1 \\
            \end{pmatrix}\right\} \\
            W&: \{\vec f_1 = (1, 0, 0), \vec f_2 = (0, 1, 0), \vec f_3 = (0, 0, 1)\}
        \end{align*}
        Now, to determine $M$, we find $T(\vec e_i)$ for each $i$:
        \[
            T(\vec e_1) = \begin{pmatrix}
             1 \\
             0 \\
             0 \\
            \end{pmatrix}, \;
            T(\vec e_2) = \begin{pmatrix}
             1 \\
             0 \\
             0 \\
            \end{pmatrix}, \;
            T(\vec e_3) = \begin{pmatrix}
             0 \\
             1 \\
             0 \\
            \end{pmatrix}, \;
            T(\vec e_4) = \begin{pmatrix}
             0 \\
             0 \\
             1 \\
            \end{pmatrix}
        \]
        Hence,
        \[
            M = \begin{pmatrix}
             1 & 1 & 0 & 0 \\
             0 & 0 & 1 & 0 \\
             0 & 0 & 0 & 1 \\
            \end{pmatrix}
        \]
        
    \end{example}
    
    \subsubsection{Matrix multiplication}

    If $S$ is a linear map $U \to V$ and $T$ is a linear map $V \to W$, with corresponding matrices $N$ and $M$ respectively, then we want the matrix corresponding to $T \circ S$ to be their product -- $MN$.
    
    \[\begin{tikzcd}
        U & V & W
        \arrow["S", from=1-1, to=1-2]
        \arrow["{T \circ S}"', curve={height=12pt}, from=1-1, to=1-3]
        \arrow["T", from=1-2, to=1-3]
    \end{tikzcd}\]

    \lecture{Lecture 12 -- 04/11/25}

    Suppose we have bases for the relevant spaces:
    \begin{align*}
        \{\vec e_1, \dots, \vec e_n\} &\text{ a basis of } V \quad (\dim V = n)\\
        \{\vec f_1, \dots, \vec f_m\} &\text{ a basis of } W \quad (\dim W = m)\\
        \{\vec g_1, \dots, \vec g_l\} &\text{ a basis of } U \quad (\dim U = l)
    \end{align*}

    Then the matrix $L = MN$ should have coefficients given by
    \begin{align*}
        L_{ik} &= M_{ij}N_{jk} \\
        &= [\vec R_i(M)]_j [\vec C_k(N)]_j = \vec R_i(M) \cdot \vec C_k(N)
    \end{align*}

    Note that
    \[
        \begin{cases}
            M \text{ is an } m \text{ by } n \text{ matrix} \\
            N \text{ is an } n \text{ by } l \text{ matrix} \\
            L \text{ is an } m \text{ by } l \text{ matrix}
        \end{cases}
    \]

    \begin{remark}
        It's clear that, to multiply matrices, the number of columns of $M$ must be equal to the number of rows of $N$. The result $L$ will have the same number of rows as $M$ and the same number of columns as $N$.
    \end{remark}
    
    \begin{remark}
        Given this, we can now reinterpret matrix-vector multiplication as a special case of matrix multiplication, where the second `matrix' has only 1 column.
    \end{remark}
    
    We have the following properties:

    \begin{itemize}
        \item \emph{Bilinearity}. $(\lambda M + \mu N)L = \lambda(ML) + \mu(NL)$, and $L(\lambda M + \mu N) = \lambda(LM) + \mu(LN)$;
        \item \emph{Associativity}. $(MN)L = M(NL)$.
    \end{itemize}
    
    for any $M, N, L$ such that these products are defined, and scalars $\lambda$, $\mu$ from the relevant field.

    \subsubsection{Matrix inverses}

    Consider three matrices $L$, $M$ and $N$, where

    \[
        \begin{cases}
            N \text{ is an } m \text{ by } n \text{ matrix} \\
            M \text{ is an } n \text{ by } m \text{ matrix} \\
            L \text{ is an } n \text{ by } m \text{ matrix}
        \end{cases}
    \]
    
    \begin{definition}
        We say that $L$ is a \emph{left inverse} of $N$ if $LN = I$, where $I$ is the $n$ by $n$ identity matrix. We say that $M$ is a \emph{right inverse} of $N$ if $NM = I$, where $I$ is the $m$ by $m$ identity matrix.
    \end{definition}
    
    \begin{proposition}
        If $N$ is square and has left inverse $L$ and right inverse $M$, $L = M$.
    \end{proposition}
    
    \begin{proof}
        Observe
        \[
            L = L(\underbrace{NM}_I) = (\underbrace{LN}_I)M = M
        \]
    \end{proof}
    This justifies denoting $L = M$ by $N^{-1}$.

    \begin{remark}
        Furthermore, if $N$ has an inverse, then $N$ is a square matrix. But not all square matrices have inverses.
    \end{remark}
    
    \begin{definition}
        A matrix that has an inverse is called \emph{invertible} or \emph{non-singular}.
    \end{definition}
    
    \begin{proposition}
        If $M$ and $N$ are invertible matrices,
        \[
            (MN)^{-1} = N^{-1}M^{-1}
        \]
    \end{proposition}
    
    \begin{proof}
        \[
            (N^{-1}M^{-1})MN = N^{-1}(M^{-1}M)N = N^{-1}N = I
        \]
    \end{proof}
    
    \begin{example}[Inverses of particular matrices]
        \begin{itemize}
            \item \emph{Rotations}. $\Rot(\theta, \hvec n)^{-1} = \Rot(-\theta, \hvec n)$.
            \item \emph{Shears}. For fixed $\vec a$, $\vec b$, $S(\lambda)^{-1} = S(-\lambda)$.
            \item \emph{Reflections}. For any reflection $H$, $H^{-1} = H$.
        \end{itemize}
    \end{example}
    
    \subsubsection{Transpose and Hermitian conjugate}

    \begin{definition}
        If $M$ is an $m$ by $n$ matrix, the \emph{transpose} $M^\intercal$ is an $n$ by $m$ matrix given by
        \[
            (M^\intercal)_{ij} = M_{ji}
        \]
        or
        \[
            \begin{pmatrix}
             M_{11} & \cdots & M_{1n} \\
             \vdots &  & \vdots \\
             M_{m1} & \cdots & M_{mn} \\
            \end{pmatrix}^\intercal =
            \begin{pmatrix}
             M_{11} & \cdots & M_{m1} \\
             \vdots &  & \vdots \\
             M_{1n} & \cdots & M_{mn} \\
            \end{pmatrix}
        \]
    \end{definition}
    
    Properties of the transpose include:

    \begin{itemize}
        \item $(M^\intercal)^\intercal = M$; 
        \item If $\vec x$ is a column vector, then $\vec x^\intercal$ is the corresponding row vector;
        \item $(NM)^\intercal = M^\intercal N^\intercal$;
        \item $(\alpha M + \beta N)^\intercal = \alpha M^\intercal + \beta N^\intercal$.
    \end{itemize}
    
    \begin{definition}
        If $M$ is a square matrix, then $M$ is
        \begin{itemize}
            \item \emph{symmetric} if $M^\intercal = M$, that is $M_{ij} = M_{ji}$ for all $i$, $j$;
            \item \emph{antisymmetric} if $M^\intercal = -M$.
        \end{itemize}
    \end{definition}
    
    \begin{definition}
        The \emph{Hermitian conjugate} of a $m$ by $n$ matrix $M$ is given by
        \[
            M^\dagger = \overline{M^\intercal}
        \]
        which is an $n$ by $m$ matrix. Clearly $(M^\dagger)_{ij} = \overline{M_{ji}}$.
    \end{definition}
    
    Properties of the Hermitian conjugate:

    \begin{itemize}
        \item $(M^\dagger)^\dagger = M$; 
        \item $(NM)^\dagger = M^\dagger N^\dagger$;
        \item $(\alpha M + \beta N)^\dagger = \bar\alpha M^\dagger + \bar\beta N^\dagger$.
    \end{itemize}
    
    \begin{definition}
        If $M$ is a square matrix, then $M$ is
        \begin{itemize}
            \item \emph{Hermitian} if $M^\dagger = M$;
            \item \emph{anti-Hermitian/skew-Hermitian} if $M^\dagger = -M$.
        \end{itemize}
    \end{definition}
    
    \subsubsection{Trace}

    \begin{definition}
        For an $n$ by $n$ matrix $M$, the \emph{trace}
        \[
            \Tr(M) = M_{ii} = \sum_{i=1}^n M_{ii}
        \]
        is the sum of the diagonal entries. $M$ is \emph{traceless} if $\Tr(M) = 0$.
    \end{definition}
    
    Properties of the trace:

    \begin{itemize}
        \item \emph{Linearity}. $\Tr(\alpha M + \beta N) = \alpha \Tr(M) + \beta \Tr(N)$.
        \item \emph{Cyclicity}. $\Tr(MN) = \Tr(NM) \neq \Tr(M)\Tr(N)$.
        \item $\Tr(M^\intercal) = \Tr(M)$.
        \item $\Tr(I) = n$ where $I$ is the $n$ by $n$ identity matrix.
    \end{itemize}
    
    \subsubsection{Decompose of an $n$ by $n$ matrix}

    Any (real) square matrix is a sum of symmetric and antisymmetric parts. That is, for such a matrix $M$, we have
    \[
        M = S + A
    \]
    where
    \[
        \begin{cases}
            S = \frac{1}{2} (M + M^\intercal) \\
            A = \frac{1}{2} (M - M^\intercal)
        \end{cases}
    \]
    and it is clear that $S$ is symmetric and $A$ is antisymmetric. The symmetric part can be further decomposed as
    \[
        S = T + \frac 1 n \Tr(S) I
    \]
    where $\Tr(T) = 0$. Furthermore, $\Tr(S) = \Tr(M)$ and $\Tr(A) = 0$. Combining these decompositions,
    \[
        M = T + A + \frac 1 n \Tr(M) I
    \]
    where $T$ is symmetric and traceless, $A$ is antisymmetric, and $\frac 1 n \Tr(M) I$ is the \emph{isotropic part} -- it is diagonal with all entries equal.

    \lecture{Lecture 13 -- 06/11/25}
    
    \subsection{Orthogonal and unitary matrices}

    \begin{definition}
        A real $n$ by $n$ matrix $U$ is \emph{orthogonal} if and only if
        \[
            UU^\intercal = U^\intercal U = I
        \]
    \end{definition}
    Equivalent properties to being orthogonal:
    \begin{itemize}
        \item $U^\intercal = U^{-1}$;
        \item the columns and rows of $U$ are orthonormal -- $U_ijU_jk = \delta_{ik}$;
        \item $U$ preserves inner products: $(U \vec x) \cdot (U \vec y) = \vec x \cdot \vec y$ for any $\vec x$, $\vec y$, and so $U$ also preserves norms and angles.
    \end{itemize}
    
    \begin{definition}
        A complex $n$ by $n$ matrix $U$ is \emph{unitary} if
        \[
            UU^\dagger = U^\dagger U = I
        \]
    \end{definition}
    Equivalently,
    \begin{itemize}
        \item $U^\dagger = U^{-1}$;
        \item $U$ preserves inner products: $(U \vec z)^\dagger (U \vec w) = \vec z^\dagger \vec w$ (we're using vector transpose notation for the inner product).
    \end{itemize}
    
    \begin{lemma}
        Orthogonal matrices in $\RR^2$ are exactly rotation or reflection matrices.
    \end{lemma}
    
    \begin{proof}
        Let $U$ be orthogonal, and consider the image of the canonical basis elements $\vec e_1$ and $\vec e_2$.
        \begin{itemize}
            \item Since $U$ preserves norms,
            \[
                U \vec e_1 = 
                U \begin{pmatrix}
                 1 \\
                 0 \\
                \end{pmatrix} = \begin{pmatrix}
                 \cos \theta \\
                 \sin \theta \\
                \end{pmatrix}
            \]
            for any $\theta$, since this is the general form of a vector of norm 1.
            \item Since $U$ preserves angles, we must have a 90 degree angle between the images of $\vec e_1$ and $\vec e_2$, so
            \[
                U \vec e_2 =
                U \begin{pmatrix}
                 0 \\
                 1 \\
                \end{pmatrix} = \pm \begin{pmatrix}
                 -\sin \theta \\
                 \cos \theta \\
                \end{pmatrix}
            \]
        \end{itemize}
        Thus,
        \begin{align*}
            U = \begin{pmatrix}
                \cos \theta & -\sin \theta \\
                \sin \theta & \cos \theta \\
            \end{pmatrix}
            = \Rot(\theta)
        \end{align*}
        or
        \begin{align*}
            U = \begin{pmatrix}
                \cos \theta & \sin \theta \\
                \sin \theta & -\cos \theta \\
            \end{pmatrix}
            = \Rfl(\theta)
        \end{align*}
    \end{proof}
    
    \subsection{The determinant}
    
    Consider an invertible map from $\RR^n \to \RR^n$ given by a real $n$ by $n$ matrix $M$.

    \subsubsection{$\RR^2$}

    Let
    \begin{align*}
        M = \begin{pmatrix}
         M_{11} & M_{12} \\
         M_{21} & M_{22} \\
        \end{pmatrix}
    \end{align*}
    and let
    \begin{align*}
        \tilde{M} = \begin{pmatrix}
         M_{22} & -M{12} \\
         -M{21} & M_{22} \\
        \end{pmatrix}
    \end{align*}
    Then,
    \begin{align*}
        \vec x' = M\vec x \implies \tilde M \vec x' &= \tilde MM \vec x \\
        &= (\det M) \vec x
    \end{align*}
    where $\det M = M_{11}M_{22} - M_{12}M_{21}$. Recalling our second scalar product in $\RR^2$, $[\vec a, \vec b] = a_1b_2 - a_2b_1$, we can recognise that $\det M = [M\vec e_1, M\vec e_2]$, and this leads to the interpretation of $\det M$ as the signed area of the parallelogram spanned by the image of the canonical basis under $M$.
    
    Anyway, if $\det M \neq 0$, we have
    \[
        M^{-1} = \frac{\tilde M}{\det M}
    \]
    and so $M$ is invertible.

    \subsubsection{$\RR^3$}

    Let us attempt once again to find such an $\tilde M$. Recalling that the scalar triple product was the 3D analogue of our second 2D inner product, we might wish to consider
    \[
        [\vec a, \vec b, \vec c] = \vec a \cdot (\vec b \times \vec c) = \lv{ijk}a_ib_jc_k
    \]
    We then define the determinant to be the signed volume of the parallelepiped spanned by the image of the canonical basis:
    \begin{align*}
        \det M = [M\vec e_1, M\vec e_2, M \vec e_3] &= [\vec C_1(M), \vec C_2(M), \vec C_3(M)] \\
        &= [M_{i1}\vec e_i, M_{j2}\vec e_j, M_{k3}\vec e_k] \\
        &= M_{i1}M_{j2}M_{k3} [\vec e_i, \vec e_j, \vec e_k] \\
        &= M_{i1}M_{j2}M_{k3} \lv{ijk}
    \end{align*}
    which we denote by
    \[
        \det M = \begin{vmatrix}
         M_{11} & M_{12} & M_{13} \\
         M_{21} & M_{22} & M_{23} \\
         M_{31} & M_{32} & M_{33} \\
        \end{vmatrix}
    \]
    Now we can actually go on to construct $\tilde M$ as we originally wanted. Let us define
    \begin{align*}
        \vec R_1(\tilde M) &= \vec C_2 (M) \times \vec C_3 (M) \\
        \vec R_2(\tilde M) &= \vec C_3 (M) \times \vec C_1 (M) \\
        \vec R_3(\tilde M) &= \vec C_1 (M) \times \vec C_3 (M)
    \end{align*}
    from which we can see
    \begin{align*}
        \vec R_i(\tilde M) \cdot \vec C_j(M) &= \delta_{ij} \det M \\
        \implies \tilde MM &= (\det M) I
    \end{align*}
    From this it follows that $\det M \neq 0 \iff \{\vec C_1(M), \vec C_2(M), \vec C_3(M)\}$ are linearly independent, which is equivalent to the image of $M$ being $\RR^3$ or the rank of $M$ being 3.
    
    \begin{remark*}
        General 3x3 determinants can also be found in terms of 2x2 determinants:
        \begin{align*}
            \begin{vmatrix}
            M_{11} & M_{12} & M_{13} \\
            M_{21} & M_{22} & M_{23} \\
            M_{31} & M_{32} & M_{33} \\
            \end{vmatrix} = M_{11} \begin{vmatrix}
             M_{22} & M_{23} \\
             M_{32} & M{33} \\
            \end{vmatrix} - M_{12} \begin{vmatrix}
             M_{21} & M_{23} \\
             M_{31} & M_{33} \\
            \end{vmatrix} + M_{13} \begin{vmatrix}
             M_{21} & M_{22} \\
             M_{31} & M_{32} \\
            \end{vmatrix}
        \end{align*}
    \end{remark*}
    
    \subsubsection{Permutations}

    We want to define a general Levi-Civita epsilon.

    \begin{definition}
        A \emph{permutation} of a set $S$ is a bijection $\varepsilon: S \to S$.
    \end{definition}
    
    \begin{definition}
        Let $S_n$ be the set of all permutations of $\{1, \dots, n\}$. Then $|S_n| = n!$. We denote $\rho \in S_n$ by
        \[
            \rho = \begin{pmatrix}
             1 & 2 & \cdots & n \\
             \rho(1) & \rho(2) & \cdots & \rho(n) \\
            \end{pmatrix}
        \]
    \end{definition}
    
    \begin{definition}
        A \emph{fixed point} of $\rho$ is $k$ such that $\rho(k) = k$. We usually omit fixed points when writing down permutations.
    \end{definition}
    
    \begin{definition}
        Two permutations are \emph{disjoint} if the numbers fixed by one are moved by the other and vice versa.
    \end{definition}
    
    \begin{example}
        \begin{align*}
            \begin{pmatrix}
             1 & 2 & 3 & 4 & 5 & 6 \\
             5 & 6 & 3 & 1 & 4 & 2\\
            \end{pmatrix} &= 
            \begin{pmatrix}
             1 & 2 & 4 & 5 & 6 \\
             5 & 6 & 1 & 4 & 2\\
            \end{pmatrix} \\
            &= \begin{pmatrix}
             1 & 4 & 5 \\
             5 & 1 & 4 \\
            \end{pmatrix} \begin{pmatrix}
             2 & 6 \\
             6 & 2 \\
            \end{pmatrix} \\
            &= (5\;4\;1)(6\;2)
        \end{align*}
        At the bottom, we see an alternative way of expressing our permutation -- in \emph{cycle notation}, where each element in the cycle is mapped to the next one, with the final one being mapped to the first.        
    \end{example}
    
    \begin{definition}
        A \emph{2-cycle} or a \emph{transposition} is a cycle of length 2.
    \end{definition}
    
    \begin{proposition}
        Any cycle can be written as a product of 2-cycles.
    \end{proposition}
    
    \begin{proof}
        Observe $(a_1\;a_2\;\cdots\;a_k) = (a_1\;a_2)(a_2\;a_3)\cdots(a_{k-1}\;a_k)$.
    \end{proof}
    
    \begin{corollary}
        Any permutation can be written as a product of 2-cycles.
    \end{corollary}
    
    \begin{proof}
        Decompose the permutation into disjoint cycles, and then express each cycle as a product of 2-cycles.
    \end{proof}
    
    \begin{definition}
        The \emph{sign} of a permutation $\varepsilon(\rho) = (-1)^r$, where $r$ is the number of 2-cycles of $\rho$ when it is written as a product of 2-cycles.

        If $\varepsilon(\rho) = 1$, we say it is an even permutation, odd otherwise.
    \end{definition}

    \begin{remark*}
        $\rho(\sigma \tau) = \rho(\sigma)\rho(\tau)$; $\rho(\sigma^{-1}) = \rho(\sigma)$.
    \end{remark*}
    
    
    
    \begin{definition}[Generalised Levi-Civita epsilon]
        \[
            \lv{j_1\cdots j_n} = \begin{cases}
                +1 & (j_1, \dots, j_n) \text{ is an even permutation of } (1, \dots, n) \\
                -1 & (j_1, \dots, j_n) \text{ is an odd permutation of } (1, \dots, n) \\
                0 & \text{o/w}
            \end{cases}
        \]
    \end{definition}
    
    \lecture{Lecture 14 -- 08/11/25}

    \subsubsection{Alternating forms}

    \begin{definition}
        For vectors $\vec v_1, \dots, \vec v_n$ in $\RR^n$ or $\CC^n$, the rank $n$ \emph{alternating form} is given by
        \begin{align*}
            [\vec v_1, \dots, \vec v_n] &= \lv{j_1\dots j_n}(\vec v_1)_{j_1}\dots (\vec v_n)_{j_n} \\
            &= \sum_{\rho \in S_n} \varepsilon(\rho) (\vec v_1)_{\rho(1)}\dots (\vec v_n)_{\rho(n)}
        \end{align*}
        Properties of alternating forms:
        \begin{itemize}
            \item \emph{Multilinear} (linear in any argument):
            \begin{align*}
                &\;[\vec v_1, \dots, \vec v_{p-1}, \alpha \vec u + \beta \vec w, \vec v_{p+1}, \dots, \vec v_n] \\
                =&\;\alpha[\vec v_1, \dots, \vec v_{p-1}, \vec u, \vec v_{p+1}, \dots, \vec v_n] + \\
                &\;\beta[\vec v_1, \dots, \vec v_{p-1}, \vec w, \vec v_{p+1}, \dots, \vec v_n]
            \end{align*}
            \item \emph{Totally antisymmetric}:
            \[
                [\vec v_{\rho(1)},\dots, \vec v_{\rho(n)}] = \varepsilon(\rho) [\vec v_1, \dots, \vec v_n]
            \]
            \item $[\vec e_1, \dots, \vec e_n] = +1$.
        \end{itemize}
    \end{definition}
    
    These three properties are actually sufficient to define alternating forms. Note that exchanging any two different vectors changes the sign of the form. Let's consider some further properties.

    \begin{itemize}
        \item If $\vec v_p = \vec v_q$ for $p \neq q$, then
        \[
            [\vec v_1, \dots, \vec v_p, \dots, \vec v_q, \dots, \vec v_n] = 0
        \]
        (this follows from antisymmetry).
        \item If
        \[
            \vec v_p = \sum_{i \neq p} \lambda_i \vec v_i
        \]
        then
        \[
            [\vec v_1, \dots, \vec v_p, \dots, \vec v_n] = 0
        \]
        (this follows from multilinearity and antisymmetry).
    \end{itemize}
    
    \begin{proposition}
        $[\vec v_1, \dots, \vec v_n] \neq 0 \iff \vec v_1, \dots, \vec v_n$ are linearly independent.
    \end{proposition}
    
    \begin{proof}
        \begin{itemize}
            \item[$\implies$] Suppose $\vec v_1, \dots, \vec v_n$ are linearly dependent. Then some $\vec v_p$ is a linear combination of the other, so $[\vec v_1, \dots, \vec v_n] = 0$ by the above property.
            \item[$\impliedby$] Suppose $\vec v_1, \dots, \vec v_n$ are linearly independent. Then, they span $\RR^n$/$\CC^n$ as relevant. In particular,
            \[
                \vec e_j = u_{ij} \vec v_i
            \]
            for suitable $u_{ij}$. It follows that
            \begin{align*}
                1 = [\vec e_1, \dots, \vec e_n] &= u_{i_11}\dots u_{i_nn} [\vec v_{i_1}\dots \vec v_{i_n}] \\
                &= u_{i_11}\dots u_{i_nn} \lv{i_1\dots i_n}[\vec v_1, \dots, \vec v_n]
            \end{align*}
            and so clearly $[\vec v_1, \dots, \vec v_n] \neq 0$.
        \end{itemize}
    \end{proof}
    
    \subsubsection{Determinants in $\RR^n$ and $\CC^n$}

    \begin{definition}
        Let $M$ be an $n$ by $n$ matrix with columns
        \[
            \vec C_i = M \vec e_i = M_{ij} \vec e_j
        \]
        Then, the \emph{determinant} of $M$, $\det M$, is given by
        \begin{align*}
            \det M &= [\vec C_1, \dots, \vec C_n] \\
            &= [M \vec e_1, \dots, M \vec e_n] \\
            &= \lv{j_1\dots j_n}M_{j_11}\dots M_{j_nn} \\
            &= \sum_{\rho \in S_n} \varepsilon(\rho) M_{\rho(1)1}\cdots M_{\rho(n)n}
        \end{align*}
        We may also denote it by
        \[
            \begin{vmatrix}
             M_{11} & \cdots & M_{1n} \\
             \vdots &  & \vdots \\
             M_{n1} & \cdots & M_{nn} \\
            \end{vmatrix}
        \]
    \end{definition}
    
    Properties of the determinant include:

    \begin{itemize}
        \item It is a \emph{multilinear} function of the columns. In particular, $\det(\lambda M) = \lambda^n \det M$ for an $n$ by $n$ matrix $M$.
        \item It is a \emph{totally antisymmetric} function of the columns. In particular, interchanging two columns changes the sign of the determinant.
        \item $\det I = 1$ for the identity matrix of any size.
        \item If two rows or columns of $M$ are identical, $\det M = 0$.
        \item If two rows or columns of $M$ are linearly dependent, $\det M = 0$.
        \item $\det M \neq 0$ if and only if the columns/rows are linearly independent of each other.
    \end{itemize}
    
    \begin{corollary}
        Under a column operation $\vec C_i \mapsto \vec C_i + \lambda \vec C_j$, or a row operation $\vec R_i \mapsto \vec R_i + \lambda \vec R_j$, $i \neq j$, $\det M$ is unchanged.
    \end{corollary}
    
    \begin{proposition}
        For a square matrix $M$,
        \[
            \det M = \det M^\intercal
        \]
    \end{proposition}
    
    \begin{proof}
        Consider a single term of the determinant $\varepsilon(\rho) M_{\rho(1)1}\dots M_{\rho(n)n}$, and let $\sigma = \rho^{-1}$. Then
        \begin{align*}
            \varepsilon(\rho)M_{\rho(1)1}\dots M_{\rho(n)n} = \varepsilon(\rho)M_{\sigma(\rho(1))\sigma(1)}\dots M_{\sigma(\rho(n))\sigma(n)}
        \end{align*}
        since this just amounts to a rearrangement of terms. But 
        \[
            \varepsilon(\rho)M_{\sigma(\rho(1))\sigma(1)}\dots M_{\sigma(\rho(n))\sigma(n)} = \varepsilon(\sigma)M_{1\sigma(1)}\dots M_{n\sigma(n)}
        \]
        since $\sigma = \rho^{-1}$. The result follows by summing over all such terms.
    \end{proof}
    
    \begin{remark*}
        This justifies us referring to the rows as well as the columns in our properties of determinants above, since the statements holding for the columns of $M^\intercal$ implies they hold for the rows of $M$.
    \end{remark*}
    
    \begin{proposition}
        For matrices $M$ and $N$, $\det(MN) = \det M \det N$.
    \end{proposition}
    
    \begin{corollary}
        If $M$ is orthogonal, $\det M = \pm 1$.
    \end{corollary}
    
    \begin{proof}
        $MM^\intercal = I$ implies $\det(MM^\intercal) = 1$. But $\det(MM^\intercal) = \det M \det M^\intercal = (\det M)^2$, and the result follows.
    \end{proof}
    
    \begin{corollary}
        If $M$ is unitary, $|\det M| = 1$.
    \end{corollary}

    \lecture{Lecture 15 -- 11/11/25}
    
    \begin{proof}[Proof of proposition]
        Let us first note that, according to the definition of the determinant in terms of signs of permutations, swapping columns an even or odd number of times gives a sign change of $\pm 1$ respectively. Now
        \begin{align*}
            \det(MN) &= \sum_\sigma \varepsilon(\sigma) (MN)_{\sigma(1)1} \dots (MN)_{\sigma(n)n} \\
            &= \sum_\sigma \varepsilon(\sigma) \sum_{k_1, \dots, k_n = 1}^n M_{\sigma(1)k_1}N_{k_11}\dots M_{\sigma(n)k_n}N_{k_nn} \\
            &= \sum_{k_1, \dots, k_n = 1}^n N_{k_11}\dots N_{k_nn} \underbrace{\sum_\sigma \varepsilon(\sigma) M_{\sigma(1)k_1}\dots M_{\sigma(n)k_n}}_{S}
        \end{align*}
        If, in $S$, $k_i = k_j$ for some $i \neq j$, then $S$ is the determinant of a matrix with two equal columns, and hence is zero. We may then restrict our attention to only those terms with all the $k_i$'s different -- in other words, those for which the $k_i$'s are a permutation of $1, \dots, n$, say $\rho$. Thus
        \begin{align*}
            \det(MN) &= \sum_\rho N_{\rho(1)1} \dots N_{\rho(n)n} \underbrace{\sum_\sigma \varepsilon(\sigma) M_{\sigma(1)\rho(1)} \dots M_{\sigma(n)\rho(n)}}_{\varepsilon(\rho) \det M} \\
            &= \det M \sum_\rho \varepsilon(\rho) N_{\rho(1) 1} \dots N_{\rho(n) n} \\
            &= \det M \det N
        \end{align*}
        as required.
    \end{proof}
    
    \begin{corollary}
        $1 = \det I = \det (MM^{-1}) = \det M \det M^{-1}$. Hence $\det M^{-1}$ = $(\det M)^{-1}$.
    \end{corollary}
    
    \subsubsection{Minors and cofactors}

    \begin{definition}
        For an $n$ by $n$ matrix $M$, consider the $n-1$ by $n-1$ matrix in wich row $i$ and column $j$ of $M$ have been removed. The determinant of this matrix, denoted $M^{ij}$, is called a \emph{minor}.
    \end{definition}
    
    \begin{definition}
        Let $M$ be an $n$ by $n$ matrix with columns $\vec C_j$ and rows $\vec R_i$. Then
        \begin{align*}
            \det M = [\vec C_1, \dots, \vec C_n] &= \sum_{i} M_{ij} \Delta_{ij} \\
            &= \sum_j M_{ij} \Delta_{ij}
        \end{align*}
        where the $\Delta_{ij}$ are called \emph{cofactors}. The cofactors are explicitly given by
        \begin{align*}
            \Delta_{ij} &= [\vec C_1, \dots, \vec C_{j-1}, \vec e_j, \vec C_{j+1}, \dots, \vec C_n] \\
            &= [\vec R_1, \dots, \vec R_{i-1}, \vec e_j, \vec R_{j+1}, \dots, \vec R_n]
        \end{align*}
        In particular, the cofactor is the determinant of the matrix obtained by replacing the entry $M_{ij}$ by 1 and all other entries in row $i$ and column $j$ by 0.
    \end{definition}
    
    \begin{remark}
        $\Delta_{ij} = (-1)^{i+j}M^{ij}$ -- we can see this by shifting the 1 to the top left, which will multiply by $(-1)^{i+j}$, and observing that the determinant of this matrix is just $M^{ij}$.
    \end{remark}
    Now we get a new formula for the determinant!
    \begin{align*}
        \det M &= \sum_i M_{ij} \Delta_{ij} = \sum_i (-1)^{i+j} M_{ij} M^{ij} \\
        &= \sum_j M_{ij} \Delta_{ij} = \sum_j (-1)^{i+j} M_{ij} M^{ij}
    \end{align*}
    where the summation convention does \emph{not} apply (the first and second lines hold for any $j$, $i$ respectively).

    \begin{theorem}[Laplace expansion formula]
        Let $M$ be an $n$ by $n$ matrix. For any $j$,
        \[
            \det M = \sum_{i=1}^n M_{ij} \Delta_{ij}
        \]
    \end{theorem}
    
    \begin{proof}
        We will introduce the bar notation to indicate part of a list or product which is missed out. So, for example, $\{1, 2, 3, \bar 4, 5\} = \{1, 2, 3, 5\}$. Then
        \begin{align*}
            \det M &= \sum_{i_1, \dots, i_n = 1}^n \lv{i_1 \dots i_n}M_{i_11}\dots M_{i_nn} \\
            &= \sum_{i_j=1}^n M_{i_jj} \sum_{i_1, \dots, \overline{i_j}, \dots, i_n=1}^n \lv{i_1 \dots i_n} M_{i_11} \dots \overline{M_{i_jj}} \dots M_{i_nn}
        \end{align*}
        Consider now $\sigma \in S_n$ the permutation that takes $i_j$ to the $j^\text{th}$ position, and leaves everything else in its original order:
        \begin{align*}
            \sigma &= \begin{pmatrix}
             1 & \dots & j-1 & j & j+1 & j+2 & \dots & i_j-1 & i_j & i_j+1 & \dots & n \\
             1 & \dots & j-1 & i_j & j & j+1 & \dots & i_j-2 & i_j-1 & i_j+1 & \dots & n
            \end{pmatrix} \\
            k &\mapsto \begin{cases}
                k & 1 \leq k < j \\
                i_k & k = j \\
                k-1 & j < k \leq i_j \\
                k & i_j < k \leq n
            \end{cases}
        \end{align*}
        (in the case $i_j > j$, analogously otherwise). Since we have to perform $|j - i_j|$ transpositions to make this permutation, $\varepsilon(\sigma) = (-1)^{j - i_j}$. Consider further the permutation $\rho \in S$ given by
        \begin{align*}
            \rho &= \begin{pmatrix}
             1 & \dots & \dots & \overline{i_j} & \dots & n \\
             1 & \dots & \overline{i_j} & \dots & \dots & n 
            \end{pmatrix} \\
            &= \begin{pmatrix}
             1 & \dots & j-1 & j & \dots & i_j - 1 & i_j + 1 & \dots & n \\
             1 & \dots & i_{j-1} & i_{j+1} & \dots & i_{i_j} & i_{i_j+1} & \dots & i_n
            \end{pmatrix} \\
            k &\mapsto \begin{cases}
                i_k & 1 \leq k < j \\
                i_{k+1} & j \leq k < i_j \\
                k & k = i_j \\
                i_k & i_j < k \leq n
            \end{cases}
        \end{align*}
        and we can see that
        \[
            \rho \sigma = \begin{pmatrix}
             1 & \dots & n \\
             i_1 & \dots & i_n \\
            \end{pmatrix}
        \]
        from which it follows that
        \begin{align*}
            \varepsilon(\rho\sigma) = \lv{i_1 \dots i_n} &= \varepsilon(\rho)\varepsilon(\sigma) \\
            &= (-1)^{j-i_j}\lv{i_1 \dots \overline{i_j} \dots i_n}
        \end{align*}
        Hence we can rewrite:
        \begin{align*}
            \det M &= \sum_{i_j=1}^n M_{i_jj} \sum_{i_1, \dots, \overline{i_j}, \dots, i_n=1}^n (-1)^{j-i_j} \lv{i_1 \dots \overline{i_j} \dots i_n} M_{i_11} \dots \overline{M_{i_jj}} \dots M_{i_nn} \\
            &= \sum_{i_j=1}^n M_{i_jj}(-1)^{j-i_j}M^{i_jj} \\
            &= \sum_{i_j=1}^n M_{i_jj} \Delta_{ij} \\
            &= \sum_{i=1}^n M_{ij} \Delta_{ij}
        \end{align*}
    \end{proof}

    \begin{remark*}
        The `technical' part of this argument which makes the proof so involved is really just rigorously justifying an expression for $\lv{i_1 \dots i_n}$ in terms of $\lv{i_1 \dots \overline{i_j} \dots i_n}$. The proof is not theoretically difficult.
    \end{remark*}
    
    Reasoning as above, since $\vec C_k = \sum_i M_{ik} \vec e_i$, then
    \begin{align*}
        [\vec C_1, \dots, \vec C_{j-1}, \vec C_k, \vec C_{j+1}, \dots, \vec C_n] &= \sum_i M_{ik} \Delta_{ij} \\
        &= \begin{cases}
            \det M & j = k \\
            0 & j \neq k
        \end{cases}
    \end{align*}
    Hence
    \begin{align*}
        \sum_i M_{ik} \Delta_{ij} &= \delta_{jk} \det M \\
        \sum_j M_{lj} \Delta_{ij} &= \delta_{li} \det M
    \end{align*}

    \lecture{Lecture 16 -- 13/11/25}

    \begin{definition}
        The \emph{adjugate} of a matrix is defined as
        \[
            \tilde M = \adj M = \Delta^\intercal
        \]
        where $\Delta$ is the matrix of cofactors.
    \end{definition}

    \begin{proposition}
        From the expression above, we see that
        \[
            \tilde MM = M\tilde M = (\det M) I
        \]
        and, if $\det M \neq 0$, then
        \[
            M^{-1} = \frac{\tilde M}{\det M}
        \]
    \end{proposition}
    
    \begin{example}
        Consider
        \[
            \begin{pmatrix}
             1 & x & 1 \\
             1 & 1 & x \\
             x & 1 & 1 \\
            \end{pmatrix}
        \]
        for real $x$. We want to compute its determinant by applying the fact that taking $\vec C_i$ to $\vec C_i + \lambda \vec C_j$ does not change the determinant. We have
        \begin{align*}
            \det M = \det \begin{pmatrix}
             1 & x & 1 \\
             1 & 1 & x \\
             x & 1 & 1 \\
            \end{pmatrix} &= \det \begin{pmatrix}
             0 & x & 1 \\
             1-x & 1 & x \\
             x-1 & 1 & 1 \\
            \end{pmatrix} && \vec C_1 \to \vec C_1 - \vec C_3 \\
            &= \det \begin{pmatrix}
             0 & x & 1 \\
             2-2x & 0 & x-1 \\
             x-1 & 1 & 1 \\
            \end{pmatrix} && \vec R_2 \to \vec R_2 - \vec R_3 \\
            &= \det \begin{pmatrix}
             0 & x & 1 \\
             2-2x & 0 & x-1 \\
             x-1 & 1-x & 0 \\
            \end{pmatrix} && \vec R_3 \to \vec R_3 - \vec R_1 \\
            &= (x-1)^2 \det \begin{pmatrix}
             0 & x & 1 \\
             -2 & 0 & 1 \\
             1 & -1 & 0 \\
            \end{pmatrix} && \text{scaling } \vec R_2, \vec R_3 \\
            &= (x-1)^2 \det \begin{pmatrix}
             0 & x+2 & 0 \\
             -2 & 0 & 1 \\
             1 & -1 & 0 \\
            \end{pmatrix} && \vec R_1 \to \vec R_1 - \vec R_2 - 2 \vec R_3 \\
            &= (x-1)^2(x+2) \begin{pmatrix}
             0 & 1 & 0 \\
             -2 & 0 & 1 \\
             1 & -1 & 0 \\
            \end{pmatrix} && \text{scaling } \vec R_1 \\
            &= -(x-1)^2(x+2) \cdot \det \begin{pmatrix}
             -2 & 1 \\
             1 & 0 \\
            \end{pmatrix} \\
            &= (x-1)^2(x+2)
        \end{align*}
    \end{example}
    
    \subsection{Systems of linear equations}

    \subsubsection{2 by 2 matrices}

    Consider the system of equations given by
    \[
        \begin{cases}
            A_{11}x_1 + A_{12}x_2 = b_1 & (1) \\
            A_{21}x_1 + A_{22}x_2 = b_2 & (2)
        \end{cases}
    \]
    We can write this as
    \begin{align*}
        A \vec x &= \vec b \\
        \text{where } A = \begin{pmatrix}
         A_{11} & A_{12} \\
         A_{21} & A_{22} \\
        \end{pmatrix}&, \vec b = \begin{pmatrix}
         b_1 \\
         b_2 \\
        \end{pmatrix}
    \end{align*}
    Then, multiplying (1) by $A_{22}$ and (2) by $A_{12}$, and taking the difference, we have
    \[
        (A_{11}A_{22} - A_{12}A_{21})x_1 = A_{22}b_1 - A_{12}b_2
    \]
    Similarly, $(2) \times A_{11} - (1) \times A_{21}$ gives
    \[
        (A_{11}A_{22} - A_{12}A_{21})x_2 = A_{11}b_2 - A_{21}b_1
    \]
    Recognising $A_{11}A_{22} - A_{12}A_{21}$ as $\det A$, we can write this as
    \[
        \begin{pmatrix}
         x_1 \\
         x_2 \\
        \end{pmatrix} = \frac{1}{\det A} \begin{pmatrix}
         A_{22} & -A_{12} \\
         -A_{21} & A_{11} \\
        \end{pmatrix} \begin{pmatrix}
         b_1 \\
         b_2 \\
        \end{pmatrix}
    \]
    However, we can clearly recognise this matrix term as $A^{-1}$, giving us the result, as expected, that $\vec x = A^{-1} \vec b$.

    \subsubsection{General case}

    Let's move on to a system with $n$ unknowns $x_i$, written in vector-matrix form:
    \[
        A\vec x = \vec b
    \]
    with $\vec x$, $\vec b \in \RR^n$ and $A$ an $n$ by $n$ matrix. We will examine three scenarios:
    \begin{enumerate}[label=\protect\circled{\arabic*}]
        \item $\det A \neq 0$, so $A^{-1}$ exists and there is a unique solution given by $\vec x = A^{-1} \vec b$.
        \item $\det A = 0$ and $\vec b \notin \Im A$, then there are no solutions;
        \item $\det A = 0$ and $\vec b \in \Im A$, then there are infinitely many solutions. This is the most interesting case.
    \end{enumerate}
    
    \begin{proposition}
        Given a solution $\vec x_0$ to an equation $A \vec x = \vec b$, all the solutions are of the form $\vec x_0 + \vec u$, where $\vec u \in \ker A$.
    \end{proposition}
    
    \begin{proof}
        \begin{itemize}
            \item[$\impliedby$] If $\vec x = \vec x_0 + \vec u$, then $A \vec x = A (\vec x_0 + \vec u) = A \vec x_0 + A \vec u = A \vec x_0 = \vec b$.
            \item[$\implies$] Suppose $A \vec x = \vec b$, then $\vec 0 = A \vec x - A \vec x_0 = A(\vec x - \vec x_0)$, so letting $\vec u = \vec x - \vec x_0$, we have $\vec x = \vec x_0 + \vec u$ with $\vec u$ in the kernel as required.
        \end{itemize}
    \end{proof}
    
    We can also reinterpret case 1 in terms of this result -- that $\ker A = \{\vec 0\}$. In particular, we can then take a basis for the kernel, and obtain
    \[
        \vec u = \sum_{i=1}^k \lambda_i \vec u_i
    \]
    where $k$ is the size of a basis for $\ker A$, equivalently $\nullity(A)$.

    \begin{example}
        Let's now solve $A \vec x = \vec b$ with
        \[
            A = \begin{pmatrix}
             1 & x & 1 \\
             1 & 1 & x \\
             x & 1 & 1 \\
            \end{pmatrix}
        \]
        as before and $\vec b = (1, y, 1)$. We saw before that $\det A = (x-1)^2(x+2)$.
        \begin{itemize}
            \item Suppose $\det A \neq 0 \iff x \neq 1, -2$. Then $A^{-1}$ exists and we have a unique solution for $\vec x$. To actually find $A^{-1}$, let us consider cofactors, so we have
            \begin{align*}
                A^{-1} = \frac{\Delta^\intercal}{\det A}, \; \Delta &= \begin{pmatrix}
                 1-x & x^2-1 & 1-x \\
                 1-x & 1-x & x^2-1 \\
                 x^2-1 & 1-x & 1-x \\
                \end{pmatrix} \\
                \implies A^{-1} &= \frac{1}{(1-x)(x+2)} \begin{pmatrix}
                 1 & 1 & -x-1 \\
                 -x-1 & 1 & 1 \\
                 1 & -x-1 & 1 \\
                \end{pmatrix}
            \end{align*}
            Hence,
            \begin{align*}
                \vec x = A^{-1} \vec b &= \frac{1}{(1-x)(x+2)} \begin{pmatrix}
                 1 & 1 & -x-1 \\
                 -x-1 & 1 & 1 \\
                 1 & -x-1 & 1 \\
                \end{pmatrix} \begin{pmatrix}
                 1 \\
                 y \\
                 1 \\
                \end{pmatrix} \\
                &= \frac{1}{(1-x)(x+2)} \begin{pmatrix}
                 y-x \\
                 -x+y \\
                 2-xy-y \\
                \end{pmatrix}
            \end{align*}
            so we have a unique solution.
            \item Suppose $x = 1$. Then
            \begin{align*}
                A = \begin{pmatrix}
                 1 & 1 & 1 \\
                 1 & 1 & 1 \\
                 1 & 1 & 1 \\
                \end{pmatrix},\; \Im A = \Span\left\{\begin{pmatrix}
                 1 \\
                 1 \\
                 1 \\
                \end{pmatrix}\right\},\; \ker A = \Span\left\{\begin{pmatrix}
                 1 \\
                 -1 \\
                 0 \\
                \end{pmatrix}, \begin{pmatrix}
                 1 \\
                 0 \\
                 -1 \\
                \end{pmatrix}\right\}
            \end{align*}
            so we have a solution if and only if $y=1$. When $y=1$, a particular solution is given by, for example, $\vec x_0 = (1, 0, 0)$. Hence, the general solution is $\vec x = \vec x_0 + \vec u$, $u \in \ker A$ -- in other words,
            \begin{align*}
                \vec x = \begin{pmatrix}
                 1+\lambda + \mu \\
                 -\lambda \\
                 -\mu \\
                \end{pmatrix}
            \end{align*}
            for any $\lambda$, $\mu \in \RR$. Geometrically, this is a plane.
            \item The case $x=-2$ proceeds similarly and is left as an exercise to the reader.
        \end{itemize}
    \end{example}
    
    \lecture{Lecture 17 -- 15/11/25}

    \subsubsection{The homogeneous case: geometric interpretation}

    Consider the equation $A \vec u = \vec 0$. Then, if $\vec R_1, \vec R_2, \vec R_3$ are the rows of $A$, we have
    \[
        \begin{cases}
            \vec R_1 \cdot \vec u = 0 \\
            \vec R_2 \cdot \vec u = 0 \\
            \vec R_3 \cdot \vec u = 0
        \end{cases}
    \]
    Observe each of these equations describes a plane through the origin with normal vector $\vec R_1$, $\vec R_2$, $\vec R_3$ respectively. Thus, the solution of $A \vec u = \vec 0$ (which is $\ker A$) is the intersection of these planes. We split into cases.
    \begin{itemize}
        \item $\rank A = 3 \iff \nullity A = 0$, so $\ker A = \{\vec 0\}$. Here the normal vectors are linearly independent and the planes intersect at $\vec 0$.
        \item $\rank A = 2 \iff \nullity A = 1$. Here the planes intersect in a line, and the three normals are coplanar.
        \item $\rank A = 1 \iff \nullity A = 2$. Here the intersection is itself a plane, so the three planes are coincident, and all normals are parallel.
    \end{itemize}
    
    \subsubsection{The general case: geometric interpretation}

    In the general case $A \vec x = \vec b$, we instead have
    \[
        \begin{cases}
            \vec R_1 \cdot \vec x = b_1 \\
            \vec R_2 \cdot \vec x = b_2 \\
            \vec R_3 \cdot \vec x = b_3
        \end{cases}
    \]
    so we have planes with normals $\vec R_1$, $\vec R_2$, $\vec R_3$, but not necessarily all passing through the origin. Again, we have cases.

    \begin{itemize}
        \item $\rank A = 3 \iff \det A \neq 0$. Here the normals are linearly independent, and the planes intersect at some point, giving a unique solution.
        \item $\rank A < 3 \iff \det A = 0$. Now, whether a solution exists depends on whether $\vec b \in \Im A$.
        \begin{itemize}
            \item $\rank A = 2$. Then the planes may intersect in a line, as in the homogeneous case, or there may be no solution.
            DIAGRAM: range of possible arrangements (prism, parallel crossed, intersection)
            \item $\rank A = 1$. Then the planes may coincide, as in the homogeneous case, or there may be no solution.
            DIAGRAM: range of possible arrangements
        \end{itemize}
    \end{itemize}
    
    \subsubsection{Gaussian elimination}

    Consider a system of $m$ equations in $n$ unknowns:

    \[
        \begin{cases}
            A_{11}x_1 + \dots + A_{1n}x_n = b_1 \\
            \qquad \vdots \\
            A_{m1}x_1 + \dots + A_{mn}x_n = b_m
        \end{cases}
    \]
    and we'll assume without loss of generality that $A_{11} \neq 0$. We'll perform a process to change the values of the coefficients $A$ and $b$ while not changing the system as a whole -- we'll denote the coefficients in the $i^\text{th}$ step with a superscript $(i)$.
    \begin{itemize}
        \item We subtract multiples of the first equation from all the others to make the coefficients of $x_1$ vanish:
        \begin{align*}
            A_{11}^{(1)}x_1 + A_{12}^{(1)}x_2 + \dots + A_{1n}^{(1)}x_n &= b_1^{(1)} \\
            A_{22}^{(1)}x_2 + \dots + A_{2n}^{(1)}x_n &= b_2^{(1)} \\
            \vdots \\
            A_{m2}^{(1)}x_2 + \dots + A_{mn}^{(1)}x_n &= b_m^{(1)}
        \end{align*}
        \item We repeat the above step for $A_{22}^{(1)}$ and the coefficients of $x_2$ and so on:
        \begin{align*}
            A_{11}^{(1)}x_1 + A_{12}^{(1)}x_2 + \dots + A_{1n}^{(1)}x_n &= b_1^{(1)} \\
            A_{22}^{(2)}x_2 + \dots + A_{2n}^{(2)}x_n &= b_2^{(2)} \\ 
            \vdots \\
            A_{rr}^{(r)}x_r + \dots + A_{rn}^{(r)}x_n &= b_r^{(r)} \\
            0 &= b_{r+1}^{(r)} \\
            \vdots \\
            0 &= b_m^{(r)}
        \end{align*}
        and in all these equations, we need $A_{ii}^{(i)} = 0$. We stop when all the rows are equal to 0.
    \end{itemize}
    Again, we should look at cases.
    \begin{itemize}
        \item $r = n \leq m$ and $b_i^{(r)} = 0$ for $i = r+1, \dots, m$. Then, there's a unique solution, which we get by finding $x_n$ and substituting back in, working upwards.
        \item $r < m$ and $b_i^{(r)} \neq 0$ for some $i = r+1, \dots, m$. Then, there is no solution.
        \item $r = m$. Then $x_{r+1}, \dots, x_n$ are undetermined. Given any $b_1, \dots, b_m$, we can solve for $x_1, \dots, x_m$, so there are infinitely many solutions. 
    \end{itemize}
    However, we can also take a matrix approach to Gaussian elimination. Given a matrix equation
    \[
        A \vec x = \vec b
    \]
    with $A$ an $m$ by $n$ matrix. We can re-express this in the form $M \vec x = \vec d$ where $M$ is of the form
    \[
        \begin{pmatrix}
         M_{11} & M_{12} & \cdots & M_{1r} & \cdots & M_{1n} \\
         0 & M_{22} & \cdots & M_{2r} & \cdots & M_{2n} \\
         \vdots & \vdots & \ddots & \vdots & \ddots & \vdots \\
         0 & 0 & \cdots & M_{rr} & \cdots & M_{rn} \\
         0 & 0 & 0 & 0 & 0 & 0 \\
         \vdots & \vdots & \vdots & \vdots & \vdots & \vdots \\
         0 & 0 & 0 & 0 & 0 & 0
        \end{pmatrix}
    \]
    This is called \emph{echelon form}. It has some cool properties!
    \begin{itemize}
        \item The top left $r$ by $r$ block is \emph{upper triangular}, and $M_{ii} \neq 0$ for $i = 1, \dots, r$.
        \item $r = \rank M = \rank A$.
        \item If $n = m$, $\det A = \pm \det M$, and if $n=m=r$, $\det M = M_{11}\dots M_{rr} \neq 0$, so both $A$ and $M$ are invertible.
    \end{itemize}
    
    \section{Eigenvectors and eigenvalues}

    \subsection{Introduction}

    Recall \cref{fta}, the \emph{fundamental theorem of algebra} -- every complex-valued polynomial has precisely $m$ (not necessarily distinct) complex roots, where $m$ is its degree.
    
    \begin{definition}
        The root $z = w$ of a polynomial $p(z)$ has \emph{multiplicity} $k$ if $(z-w)^k$ is a factor of $p(z)$ but $(z-w)^{k+1}$ is not.
    \end{definition}
    
    \begin{definition}
        Let $T: V \to V$ (for $V$ a real or complex vector space) a linear map. Then, a vector $\vec v \in V$ with $\vec v \neq 0$ is an \emph{eigenvector} of $T$ if
        \[
            T(\vec v) = \lambda \vec v
        \]
        for some scalar $\lambda$ called an \emph{eigenvalue}.
    \end{definition}
    
    If $V = \RR^n$ or $\CC^n$, and $T$ is given as an $n$ by $n$ matrix $A$, then
    \[
        A\vec v = \lambda \vec v \iff (A - \lambda I) \vec v = \vec 0
    \]
    and, for a given $\lambda$, this holds for some non-zero vector $\vec v \neq 0$ if and only if
    \[
        \det(A - \lambda I) = 0
    \]
    This is called the \emph{characteristic equation} of $A$.

    \begin{definition}
        The \emph{characteristic equation} or \emph{characteristic polynomial} $\chi_A(t)$ of an $n$ by $n$ matrix $A$ is given by $\det(A - tI)$.
    \end{definition}
    
    \lecture{Lecture 18 -- 18/11/25}
    
    \begin{remark*}
        From the definition of the determinant,
        \begin{align*}
            \chi_A(t) &= \det(A - t I) \\
            &= \lv{j_1\dots j_n}(A_{j_11}-t\delta_{j_11})\dots(A_{j_nn}-t\delta_{j_nn}) \\
            &= c_0 + c_1t + \dots + c_nt^n
        \end{align*}
        for some constants $c_0, \dots, c_n$.
        \begin{itemize}
            \item Hence, $\chi_A(t)$ has degree $n$, and thus $n$ roots (counting multiplicity) -- we conclude an $n$ by $n$ matrix has $n$ eigenvalues.
            \item If $A$ is real, then all $c_i \in \RR$, so the eigenvalues are all either real or come in conjugate pairs.
            \item $c_n = (-1)^n$, since we multiply $-t$ with itself $n$ times, and $c_{n-1} = (-1)^{n-1}(A_{11} + \dots + A_{nn}) = (-1)^{n-1} \Tr A$. But $c_{n-1}/c_n$ is also the sum of roots, so
            \[
                \Tr A = t_1 + \dots + t_n,
            \]
            the sum of all the eigenvalues.
            \item $c_0 = \chi_A(0) = \det A$. But $c_0/c_n$ is also the product of all the negative roots, so $c_0/(-1)^n = (-1)^n t_1 \dots t_n$, and hence
            \[
                \det A = t_1 \dots t_n
            \]
        \end{itemize}
    \end{remark*}
    
    \begin{example}
        \begin{itemize}
            \item Take $V = \CC^2$,
            \[
                A = \begin{pmatrix}
                 0 & -1 \\
                 1 & 0 \\
                \end{pmatrix}
            \]
            a 90 degree rotation. Then
            \[
                \chi_A(t) = \det(A - tI) = \det\begin{pmatrix}
                 -t & -1 \\
                 1 & -t \\
                \end{pmatrix} = t^2+1
            \]
            and so the eigenvalues are $\pm i$. For $\lambda = i$, we find the eigenvectors by solving the equation
            
            \begin{align*}
                \begin{pmatrix}
                 -i & -1 \\
                 1 & -i \\
                \end{pmatrix} \begin{pmatrix}
                 x \\
                 y \\
                \end{pmatrix} &= \begin{pmatrix}
                 0 \\
                 0 \\
                \end{pmatrix} \\
                \implies \begin{cases}
                    -ix - y = 0 \\
                    x - iy = 0
                \end{cases}    
            \end{align*}
            which is solved by any vector of the form $\alpha(1, -i)$. For $\lambda = -i$ our eigenvectors are of the form $\alpha(1, i)$.
            \item Take $V = \RR^2$,
            \[
                A = \begin{pmatrix}
                 1 & 1 \\
                 0 & 1 \\
                \end{pmatrix}
            \]
            with $\chi_A(t) = (t-1)^2$, and hence eigenvalue 1 with multiplicity 2. It can be found that the eigenvectors are of the form $\alpha(1, 0)$.
        \end{itemize}
    \end{example}
    
    \subsection{Eigenspaces and multiplicity}

    \begin{definition}
        For an eigenvalue $\lambda$ of a matrix $A$, the \emph{eigenspace} of $\lambda$ is
        \begin{align*}
            E_\lambda &= \{\vec v : A \vec v = \lambda \vec v\} \\
            &= \ker(A - \lambda I)
        \end{align*}
    \end{definition}
    
    \begin{definition}
        The \emph{algebraic multiplicity} $M(\lambda) = M_\lambda$ of an eigenvalue $\lambda$ is the multiplicity of $\lambda$ in the characteristic polynomial.
    \end{definition}
    
    \begin{remark*}
        By the fundamental theorem of algebra,
        \[
            \sum_\lambda M(\lambda) = n
        \]
        for an $n$ by $n$ matrix $A$.
    \end{remark*}
    
    \begin{definition}
        The \emph{geometric multiplicity} $m(\lambda) = m_\lambda$ of an eigenvalue $\lambda$ is the dimension of the eigenspace of $\lambda$, i.e. the maximum number of linearly independent eigenvectors with eigenvalue $\lambda$.
    \end{definition}
    
    \begin{proposition}
        Consider an $n$ by $n$ matrix $A$ and $\lambda$ an eigenvalue. We have
        \[
            M_\lambda \geq m_\lambda
        \]
        \label{geometric-algebraic-multiplicity}
    \end{proposition}
    We don't have time to prove this, sadly.
    
    \begin{definition}
        Given an eigenvalue $\lambda$, its \emph{defect} $\Delta_\lambda$ is given by
        \[
            \Delta_\lambda = M_\lambda - m_\lambda
        \]
        which is $\geq 0$ by the proposition.
    \end{definition}
    
    \begin{example}
        \begin{itemize}
            \item Consider
            \[
                A = \begin{pmatrix}
                 4 & 1 & 0 \\
                 0 & 4 & 1 \\
                 0 & 0 & 4 \\
                \end{pmatrix}
            \]
            with $\chi_A(T) = (4-t)^3$ and thus eigenvalue $4$ with algebraic multiplicity $M_4 = 3$. To find the algebraic multiplicity, we must examine
            \[
                A - 4I = \begin{pmatrix}
                 0 & 1 & 0 \\
                 0 & 0 & 1 \\
                 0 & 0 & 0 \\
                \end{pmatrix}
            \]
            whose kernel -- that is the eigenspace $E_4$ -- is spanned by $(1, 0, 0)$. Thus $m_4 = \dim(E_4) = 1$.
            \item Consider a reflection in a plane through $\vec 0$ with normal $\hvec n$, satisfying $H \hvec n = -\hvec n$, $H \vec u = \vec u$ for every $\vec u \perp \hvec n$. By exercising our capacity for imagination, we have $\lambda = \pm 1$, with
            \begin{align*}
                E_{-1} &= \Span \{\hvec n\} &&\implies m_{-1} = 1\\
                E_1 &= \Span \{\vec x : \vec x \cdot \hvec n = 0\} &&\implies m_1 = 2
            \end{align*}
            Then, since $M_{-1} \geq m_{-1}$ and $M_1 \geq M_{-1}$, we can apply $M_{-1} + M_1 = 3$ to conclude $M_{-1} = m_{-1} = 1$ and $M_1 = m_1 = 2$.
            \item Consider a rotation in $\RR^2$
            \[
                \Rot(\theta) = \begin{pmatrix}
                 \cos \theta & -\sin \theta \\
                 \sin \theta & \cos \theta \\
                \end{pmatrix}
            \]
            so $\chi_{\Rot(\theta)}(t) = t^2 - 2t \cos \theta + 1$ with eigenvalues $e^{\pm i \theta}$ and corresponding eigenvectors $\vec v = \alpha (1, \mp i)$.
            \item Consider a rotation in $\RR^3$ by angle $\theta$ about $\hvec n$, $\Rot(\theta; \hvec n)$. Again it is easy to see we have an eigenvalue of 1 with eigenspace
            \[
                E_\lambda = \Span \{\hvec n\}
            \]
            Furthermore, there are no other real eigenvalues unless $\theta = n\pi$, since we reduce to the 2D rotation case, where we already know the eigenvalues are $e^{\pm i \theta}$.
            \item Consider
            \[
                A = \begin{pmatrix}
                 -3 & -1 & 1 \\
                 -1 & -3 & 1 \\
                 -2 & -2 & 0 \\
                \end{pmatrix}
            \]
            with $\chi_A(T) = -(t+2)^3$, and hence eigenvalue $-2$ with algebraic multiplicity 3. It can be seen that $A + 2I$ has rank 1 and nullity 2, so $-2$ has geometric multiplicity 2 and hence defect 1 -- the eigenvectors do not form a basis of $\CC^3$.
        \end{itemize}
    \end{example}
    
    \subsection{Diagonalisation and similarity}

    \begin{proposition}
        For an $n$ by $n$ matrix $A$ acting on $V = \RR^n$ or $\CC^n$, the following are equivalent:
        \begin{itemize}
            \item There exists a basis of eigenvectors for $V$, $\{\vec v_1, \dots, \vec v_n\}$ with $A\vec v_i = \lambda_i \vec v_i$ for every $i$.
            \item There exists an $n$ by $n$ invertible matrix $P$ such that
            \[
                P^{-1}AP = D = \begin{pmatrix}
                 \lambda_1 & \cdots & 0 \\
                 \vdots & \ddots & \vdots \\
                 0 & \cdots & \lambda_n \\
                \end{pmatrix}
            \]
        \end{itemize}
        \label{diagonalisability-conditions}
    \end{proposition}
    
    \begin{definition}
        An $n$ by $n$ matrix $A$ is \emph{diagonalisable} if the above hold.
    \end{definition}
    
    \lecture{Lecture 19 -- 20/11/25}

    \subsubsection{Linearly independent eigenvectors}

    \begin{theorem}
        Suppose that an $n$ by $n$ matrix $A$ has distinct eigenvalues $\lambda_1, \dots, \lambda_r$ with corresponding eigenvectors $\vec v_1, \dots, \vec v_r$. Then, they are linearly independent.
    \end{theorem}
    
    \begin{proof}
        By contradiction. Suppose that $\{\vec v_1, \dots, \vec v_r\}$ are linearly dependent, so
        \[
            \sum_{j=1}^r \alpha_j \vec v_j = \vec 0
        \]
        with not all $\alpha_j = 0$. We may take the minimal $p$ (potentially reordering) such that
        \[
            \sum_{j=1}^p \alpha_j \vec v_j = \vec 0
        \]
        with $\alpha_1, \dots, \alpha_p \neq 0$. Then, applying the matrix $A - \lambda_1I$ to this expression, we have
        \begin{align*}
            (A - \lambda_1I)\sum_{j=1}^p \alpha_j \vec v_j &= \sum_{j=2}^\infty \alpha_j (\lambda_j - \lambda_1)\vec v_j = \vec 0
        \end{align*}
        But now we have another linear combination of vectors with $p-1$ non-zero coefficients (since $\lambda_i \neq \lambda_j$ for $i \neq j$). This is a contradiction.
    \end{proof}
    
    \begin{remark}
        Let $B_\lambda$ be a basis for the eigenspace $E_\lambda$ associated to $\lambda$. If $\lambda_1, \dots, \lambda_r$ are distinct, then $B_{\lambda_1} \cup \dots \cup B_{\lambda_r}$ is a linearly independent set.
    \end{remark}

    \begin{proof}[Proof of \cref{diagonalisability-conditions}]
        For any matrix $P$,
        \begin{itemize}
            \item $AP$ has columns $A\vec C_i(P)$;
            \item $PD$ has columns $\lambda_i \vec C_i(P)$.
        \end{itemize}
        Thus,
        \[
            P^{-1}AP = D \iff AP = PD \iff A \vec C_i(P) = \lambda_i \vec C_i(P)
        \]
        Thus, given a basis of eigenvectors, we get a matrix $P$. Conversely, given an invertible $P$, its columns are a basis of eigenvectors (since they must be linearly independent), and we're immediately done.
    \end{proof}
    
    \subsubsection{Criteria for diagonalisability}

    \begin{itemize}
        \item \emph{Sufficient but not necessary}. If $A$ is $n$ by $n$ with $n$ distinct eigenvalues, then we have $n$ linearly independent eigenvectors, so they form a basis and hence $A$ is diagonalisable.
        \item \emph{Necessary and sufficient}. If $m_\lambda = M_\lambda$ for every eigenvalue $\lambda$, then $A$ is diagonalisable. To see this, observe that, if $\lambda_1, \dots, \lambda_r$ are all the eigenvalues of a matrix, then $B_{\lambda_1} \cup \dots \cup B_{\lambda_r}$ forms a linearly independent set, and its number of elements is
        \[
            \sum_{i} m_{\lambda_i}
        \]
        If $m_{\lambda_i} = M_{\lambda_i}$ for every $i$, then this sum is equal to $n$ and hence the eigenvectors form a basis, meaning that $A$ is diagonalisable.
    \end{itemize}
    
    \subsubsection{Similarity}

    \begin{definition}
        Two $n$ by $n$ matrices $A$ and $B$ are \emph{similar} if
        \[
            B = P^{-1}AP
        \]
        for some invertible $n$ by $n$ matrix $P$. Intuitively, similar matrices represent the same linear map with respect to different bases.
    \end{definition}

    \begin{remark}
        A matrix that is similar to a diagonal matrix is diagonalisable.
    \end{remark}
    
    \begin{proposition}
        If $A$ and $B$ are similar, then $\chi_A(t) = \chi_B(t)$. 
    \end{proposition}

    \begin{proof}
        Left as an exercise.
    \end{proof}
    
    \begin{corollary}
        It immediately follows that $\Tr A = \Tr B$ and $\det A = \det B$.
    \end{corollary}
    
    \subsubsection{Hermitian and symmetric matrices}
    
    Recall the definitions of Hermitian and symmetric matrices, and the complex and real inner products.

    \begin{remark*}
        If $A$ is Hermitian,
        \[
            (A \vec v)^\dagger \vec w = \vec v^\dagger (A \vec w)
        \]
    \end{remark*}
    
    \begin{theorem}
        For an $n$ by $n$ Hermitian matrix $A$,
        \begin{itemize}
            \item Every eigenvalue of $A$ is real;
            \item Eigenvectors $\vec v$, $\vec w$ with distinct eigenvalues are orthogonal -- that is, $\vec v^\dagger \vec w = 0$.
            \item If $A$ is symmetric, then for each eigenvalue $\lambda$, we can choose a real eigenvector $\vec v$ so that $\vec v^\intercal \vec w = 0$.
        \end{itemize}
    \end{theorem}
    
    \begin{proof}
        \begin{itemize}
            \item Consider an eigenvector $\vec v$ with eigenvalue $\lambda$. Then
            \begin{align*}
                \vec v^\dagger (A \vec v) &= (A \vec v)^\dagger \vec v \\
                \implies \vec v^\dagger (\lambda \vec v) &= (\lambda \vec v)^\dagger \vec v \\
                \implies \lambda \vec v^\dagger \vec v &= \bar\lambda \vec v^\dagger \vec v
            \end{align*}
            Then, since $\vec v \neq \vec 0$, $\lambda = \bar\lambda \implies \lambda \in \RR$.
            \item Let $\vec v$, $\vec w$ be eigenvectors with eigenvalues $\lambda$, $\mu$. Then
            \begin{align*}
                \vec v^\dagger(A \vec w) &= (A \vec v)^\dagger \vec w \\
                \implies \vec v^\dagger(\mu \vec w) &= (\lambda \vec v)^\dagger \vec w \\
                \implies \mu \vec v^\dagger \vec w &= \bar\lambda \vec v^\dagger \vec w = \lambda \vec v^\dagger \vec w
            \end{align*}
            since $\lambda \in \RR$. But, since $\mu \neq \lambda$, we conclude $\vec v^\dagger \vec w = 0$, that is they are orthogonal as required.
            \item Given a real matrix A with eigenvalue $\lambda$ and eigenvector $\vec v$, where we know $\lambda$ real, let us put $\vec v = \vec u + i \vec u'$ for $\vec u$, $\vec u' \in \RR^n$. Then
            \begin{align*}
                A \vec u &= \lambda \vec u \\
                A \vec u' &= \lambda \vec u'
            \end{align*}
            by simply taking real and imaginary parts on both sides. Then, since $\vec v \neq 0$, at least one of $\vec u$, $\vec u' \neq 0$, giving us at least one real eigenvector.
            
        \end{itemize}
        
        
    \end{proof}
    
    \subsubsection{Gram-Schmidt orthogonalisation}

    Given a linearly independent set of vectors in $\CC^n$, $\{\vec w_1, \dots, \vec w_r\}$, we can construct a sequence of sets of the form
    \begin{itemize}
        \item $\{\vec u_1, \vec w_2', \dots, \vec w_r'\}$ \\
        \item $\{\vec u_1, \vec u_2, \dots, \vec w_r'\}$ \\[5pt]
        \vdots
        \item $\{\vec u_1, \vec u_2, \dots, \vec u_r\}$
    \end{itemize}
    such that every set has the same span, all are linearly independent, and $\vec u_i$ are orthonormal to one another and the $\vec w$'s.
    
    \begin{algorithm}[Gram-Schmidt orthogonalisation]
        \begin{itemize}
            \item Since $|\vec w_1| \neq 0$, let
            \[
                \vec u_1 = \frac{\vec w_1}{|\vec w_1|}, \; \vec w_j' = \vec w_j - (\vec u_1^\dagger \vec w_j)\vec u_1
            \]
            This guarantees $|\vec u_1| = 1$ and $\vec u_1^\dagger \vec w_j' = 0$ for $j > 1$.
            \item Let 
            \[
                \vec u_2 = \frac{\vec w_2}{|\vec w_2|}, \; \vec w_j'' = \vec w_j' - (\vec u_2^\dagger \vec w_j)\vec u_2
            \]
            This guarantees $|\vec u_2| = 1$ and
            \[
                \begin{cases}
                    \vec u_2^\dagger \cdot \vec u_1 = 0 \\
                    \vec u_1^\dagger \cdot \vec w_j'' = 0 \\
                    \vec u_2^\dagger \cdot \vec w_j'' = 0
                \end{cases}
            \]
            for $j > 2$.
            \item and so on.
        \end{itemize}
    \end{algorithm}
    
    By applying this construction up to step $r$, we can find an orthonormal basis for any subspace spanned by given vectors. In particular, we can find an orthonormal basis $B_\lambda$ for each eigenspace $E_\lambda$ of a Hermitian matrix $A$. Then, for distinct eigenvalues $\lambda_1, \dots, \lambda_r$, $B_{\lambda_1} \cup \cdots \cup B_{\lambda_r}$ is an orthonormal set (since eigenvectors with distinct eigenvalues are orthogonal).

    \subsubsection{Unitary and diagonal orthogonalisation}

    \begin{theorem}
        Let $A$ be an $n$ by $n$ Hermitian matrix. Then, $A$ is diagonalisable. In particular,
        \begin{itemize}
            \item there exists a basis of eigenvectors $\vec u_1, \dots \vec u_n \in \CC^n$;
            \item there exists an $n$ by $n$ invertible matrix $P$ with
            \[
                P^{-1}AP = \begin{pmatrix}
                 \lambda_1 & \cdots & 0 \\
                 \vdots & \ddots & \vdots \\
                 0 & \cdots & \lambda_n \\
                \end{pmatrix}
            \]
            where the columns of $P$ are the eigenvectors $\vec u_i$.
        \end{itemize}
    \end{theorem}
    
    \begin{remark*}
        In addition, the eigenvectors $\vec u_i$ may be chosen to be orthonormal, which is the same thing as $P$ being unitary.
    \end{remark*}
    
    \begin{remark*}
        If $A$ is real, then $A$ is symmetric. Then $P$ can be made to be orthogonal.
    \end{remark*}
    
    \subsection{Change of basis}

    Let $V$, $W$ be real or complex vector spaces, with $\dim V = n$, $\dim W = m$, and let $T: V \to W$ be a linear map. Fix bases $\{\vec e_1, \dots \vec e_n\}$ and $\{\vec f_1, \dots, \vec f_m\}$ of $V$ and $W$ respectively, so we may represent $T$ by the $m$ by $n$ matrix $A$ satisfying
    \[
        T(\vec e_i) = A_{ji} \vec f_i
    \]
    Fix new bases $\{\vec e_1', \dots, \vec e_n'\}$ and $\{\vec f_1', \dots, \vec f_m'\}$, so we may also represent $T$ by the $m$ by $n$ matrix $B$ satisfying
    \[
        T(\vec e_i') = B_{ji} \vec f_i'
    \]
    How are the matrices $A$ and $B$ related? Suppose the bases are related by
    \[
        \vec e_i' = P_{ki}\vec e_k, \; \vec f_i' = Q_{lj} \vec f_l
    \]
    so that $P$ and $Q$ are invertible matrices of size $n$ by $n$ and $m$ by $m$ respectively. Then

    \begin{proposition}[Change of basis]
        $B = Q^{-1}AP$.
    \end{proposition}
    
    \begin{proof}
        We have
        \begin{align*}
            T(\vec e_i') &= T(P_{ki}\vec e_k) \\
            &= P_{ki} T(\vec e_k) \\
            &= A_{jk} P_{ki} \vec f_j \\ 
        \end{align*}
        Also,
        \begin{align*}
            T(\vec e_i') &= B_{ji} \vec f_i' \\
            &= Q_{lj} B_{ji} \vec f_l \\
            &= Q_{jk} B_{ki} \vec f_j && (l \to j, j \to k)
        \end{align*}
        Hence, comparing term by term, $A_{jk}P_{ki} = Q_{jk}B_{ki}$. Hence $AP = QB$ so, multiplying on both sides by $Q^{-1}$, we have $B = Q^{-1}AP$ as required.
    \end{proof}
    
    \begin{example}
        Consider $\dim V = n = 2$, $\dim W = m = 3$, and $T$
        \begin{align*}
            \vec e_1 &\mapsto \vec f_1 + 2 \vec f_2 - \vec f_3 \\
            \vec e_2 &\mapsto -\vec f_1 + 2 \vec f_2 + \vec f_3
        \end{align*}
        so that
        \[
            A = \begin{pmatrix}
             1 & -1 \\
             2 & 2 \\ 
             -1 & 1
            \end{pmatrix}
        \]
        Now, consider a basis for $V$ $\{\vec e_1', \vec e_2'\}$, related to $\{\vec e_1, \vec e_2\}$ by $\vec e_1' = \vec e_1 - \vec e_2$, $\vec e_2' = \vec e_1 + \vec e_2$. Then
        \[
            P = \begin{pmatrix}
             1 & 1 \\
             -1 & 1 \\
            \end{pmatrix}
        \]
        For $W$, take $\{\vec f_1', \vec f_2', \vec f_3'\}$ such that $\vec f_1' = \vec f_1 - \vec f_3$, $\vec f_2' = \vec f_2$, $\vec f_3' = \vec f_1 + \vec f_3$. Then
        \[
            Q = \begin{pmatrix}
             1 & 0 & 1 \\
             0 & 1 & 0 \\
             -1 & 0 & 1 \\
            \end{pmatrix}
        \]
        Applying the change of basis formula, we get
        \[
            B = Q^{-1}AP = \begin{pmatrix}
             2 & 0 \\
             0 & 4 \\ 
             0 & 0
            \end{pmatrix}
        \]
        so we can see $T$ satisfies
        \begin{align*}
            \vec e_1' &\mapsto 2 \vec f_1' \\
            \vec e_2' &\mapsto 4 \vec f_2'
        \end{align*}
    \end{example}
    
    \begin{remark*}[Special cases]
        \begin{itemize}
            \item If $V = W$ with the same bases for both spaces, then the change of basis matrices $P$ and $Q$ are equal, so $B = P^{-1}AP$. So matrices represent the same linear map $T: V \to V$ if and only if they are similar.
            \item If $V = W = \RR^n$ or $\CC^n$, and we take the canonical basis $\{\vec e_i\}$. Then, if there exists a basis of eigenvectors of $T$, $\{\vec v_1, \dots, \vec v_n\}$, let $\vec e_i' = \vec v_i$ and let $B$ be the matrix of $T$ with respect to this basis. Then
            \[
                B = \begin{pmatrix}
                 \lambda_1 & \cdots & 0 \\
                 \vdots & \ddots & \vdots \\
                 0 & \cdots & \lambda_n \\
                \end{pmatrix} = P^{-1}AP
            \]
            Then, the columns of $P$ are the eigenvectors, and $P$ is both the change of basis matrix and the matrix that diagonalises $A$.
        \end{itemize}
    \end{remark*}

    \lecture{Lecture 21 -- 25/11/25}
    
    \subsubsection{Changes in vector components under change of basis}
    
    Let $V$ be a vector space with bases $\{\vec e_i\}$, $\{\vec e_i'\}$, and $\vec x \in V$. Let $P$ be the matrix taking $\{\vec e_i\}$ to $\{\vec e_i'\}$, and let $\{x_i\}$, $\{x_i'\}$ be such that
    \[
        \vec x = x_i \vec e_i = x_i' \vec e_i'
    \]
    Then, since $\vec e_j' = \vec e_i P_{ij}$,
    \begin{align*}
        \vec x = x_i \vec e_i &= x_j' \vec e_j' = x_j'(P_{ij} \vec e_i) = (P_{ij}x_j')\vec e_i \\
        \implies x_i &= P_{ij} x_j'
    \end{align*}
    In other words,
    \[
        \vec x = P \vec x'
    \]
    Doing the same thing with $\vec y$, $\vec y' \in W$ with bases $\{\vec f_i\}$ and $\{\vec f_i'\}$, we have $\vec y = Q\vec y'$. Taking a linear map $T: V \to W$ represented by $A$ in the normal basis and $B$ in the prime basis, we have
    \[
        B \vec x' = \vec y' = Q^{-1} \vec y = Q^{-1}A \vec x = (Q^{-1}AP) \vec x'
    \]
    and we recover the same result.

    \subsection{Cayley-Hamilton theorem}

    \begin{theorem}
        Let $A$ be an $n$ by $n$ matrix with $\chi_A(t) = \det(A - tI) = \sum_{i=0}^n c_rt^r$. Then
        \[
            \chi_A(A) = \sum_{i=0}^n c_r A^r = 0
        \]
        where $0$ is the zero matrix.
    \end{theorem}
    
    \begin{proof}
        We take cases.
        \begin{itemize}
            \item \emph{2 by 2 matrices}. One can simply check this. Letting
            \[
                A = \begin{pmatrix}
                 A_{11} & A_{12} \\
                 A_{21} & A_{22} \\
                \end{pmatrix}
            \]
            we have
            \[
                \chi_A(t) = t^2 - (A_{11}+A_{22})t + (A_{11}A_{22}-A_{12}A_{21})
            \]
            and, putting $t=A$, we find that it just works.
            \item \emph{Diagonalisable $n$ by $n$ matrices}. Consider $A$ with eigenvalues $\lambda_i$ and invertible $P$ with
            \[
                P^{-1}AP = D = \begin{pmatrix}
                 \lambda_1 & \cdots & 0 \\
                 \vdots & \ddots & \vdots \\
                 0 & \cdots & \lambda_n \\
                \end{pmatrix}
            \]
            Furthermore, we have
            \[
                D^r = \begin{pmatrix}
                 \lambda_1^r & \cdots & 0 \\
                 \vdots & \ddots & \vdots \\
                 0 & \cdots & \lambda_n^r \\
                \end{pmatrix}
            \]
            It follows that
            \[
                \chi_A(D) = \begin{pmatrix}
                 \chi_A(\lambda_1) & \cdots & 0 \\
                 \vdots & \ddots & \vdots \\
                 0 & \cdots & \chi_A(\lambda_n) \\
                \end{pmatrix} = 0
            \]
            where 0 is the zero matrix. Then, since $A = PDP^{-1}$, we have $A = PD^rP^{-1}$, and so
            \begin{align*}
                \chi_A(A) &= \chi_A(PDP^{-1}) \\
                &= P\chi_A(D)P^{-1} \\
                &= 0
            \end{align*}
            \item \emph{General case (not examinable)}. Let $M = A - tI$, so $\chi_A(t) = \det M = \sum_{r=0}^n c_rt^r$. Recall that the adjugate of $M$, $\tilde M$, satisfies $\tilde MM = (\det M)I$, and we actually know that the entries of $\tilde M$ are polynomials of degree $n-1$, so
            \[
                \tilde M = \sum_{r=0}^{n-1} B_r t^r
            \]
            Then
            \begin{align*}
                \tilde MM &= \left(\sum_{r=0}^{n-1} B_r t^r\right)(A - tI) \\
                &= B_0A + (B_1A - B_0)t + \dots + (B_{n-1}A - B_{n-2})t^{n-1} - B_{n-1}t^n
            \end{align*}
            \[
                \implies \begin{cases}
                    c_0I &= B_0A \\
                    c_1I &= B_1A-B_0 \\
                    \vdots \\
                    c_{n-1}I &= B_{n-1}A - B_{n-2} \\
                    c_nI &= -B_{n-1}
                \end{cases}
            \]
            Now, evaluating both polynomials in A, we have
            \[
                \implies \begin{cases}
                    c_0I &= B_0A \\
                    c_1A &= B_1A^2-B_0A \\
                    \vdots \\
                    c_{n-1}A^{n-1} &= B_{n-1}A^n - B_{n-2}A^{n-1} \\
                    c_nA^n &= -B_{n-1}A^n
                \end{cases}
            \]
            Summing the RHSes, we get the zero matrix, while the LHS is $\chi_A(A)$, so $\chi_A(A) = 0$ as required.
        \end{itemize}
    \end{proof}
    
    \subsection{Quadratic forms}

    \begin{definition}
        A \emph{quadratic form} is a function $\mathcal F: \RR^n \to R$ satisfying
        \[
            \mathcal F(\vec x) = \vec x^\intercal A \vec x = x_i A_{ij} x_j
        \]
        where $A$ is a real symmetric matrix.
    \end{definition}
    Then $A$ is diagonalisable, so
    \[
        P^\intercal AP = D = \begin{pmatrix}
            \lambda_1 & \cdots & 0 \\
            \vdots & \ddots & \vdots \\
            0 & \cdots & \lambda_n \\
        \end{pmatrix}
    \]
    where $\lambda_i$ are real eigenvalues and $P$ is a real orthogonal matrix of size $n$ by $n$ with columns $\vec C_i$ being orthonormal eigenvectors of $A$.

    Setting $\vec x' = P^\intercal \vec x$, $\vec x = P \vec x'$, we can diagonalise $\mathcal F$:
    \begin{align*}
        \mathcal F(\vec x) &= \vec x^\intercal A \vec x \\
        &= (P \vec x')^\intercal A (P \vec x') \\
        &= (\vec x')^\intercal P^\intercal A P \vec x' \\
        &= (\vec x')^\intercal D \vec x'
    \end{align*}
    Thus, we have
    \[
        \mathcal F(\vec x) = (\vec x')^\intercal D \vec x' = \sum_{i} \lambda_i (x_i')^2
    \]
    Furthermore, we have $x_i' = \vec C_i \cdot \vec x$ -- these are the components of $\vec x$ with respect to the new orthonormal basis vectors $\vec C_i$, which are called the \emph{principal axes} of $\mathcal F$.

    Since these are related to the standard axes by P orthogonal,
    \[
        |\vec x|^2 = x_ix_i = x'_i x'_i = |\vec x'|^2
    \]
    
    \begin{example}
        Let us work in $\RR^2$, $\mathcal F(\vec x) = \vec x^\intercal A \vec x$, where
        \[
            A = \begin{pmatrix}
                \alpha & \beta \\
                \beta & \alpha \\
            \end{pmatrix}
        \]
        Then $A$ has eigenvalues $\lambda_1 = \alpha + \beta$, $\lambda_2 = \alpha - \beta$, with corresponding eigenvectors
        \[
            \vec u_1 = \frac{1}{\sqrt 2} \begin{pmatrix}
                1 \\
                1 \\
            \end{pmatrix}, \; \vec u_2 = \frac{1}{\sqrt 2} \begin{pmatrix}
                -1 \\
                1 \\
            \end{pmatrix}
        \]
        So we have
        \begin{align*}
            \mathcal F(\vec x) = \begin{pmatrix}
             x_1 & x_2 \\
            \end{pmatrix} \begin{pmatrix}
             \alpha & \beta \\
             \beta & \alpha \\
            \end{pmatrix} \begin{pmatrix}
             x_1 \\
             x_2 \\
            \end{pmatrix} &= \alpha x_1^2 + 2\beta x_1x^2 + \alpha x_2^2 \\
            &= (\alpha + \beta)(x_1')^2 + (\alpha-\beta)(x_2')^2
        \end{align*}
        where
        \begin{align*}
            x_1 &= \vec u_1 \cdot \vec x = \frac{1}{\sqrt 2}(x_1 + x_2) \\
            x_2 &= \vec u_2 \cdot \vec x = \frac{1}{\sqrt 2}(-x_1 + x_2)
        \end{align*}
        Taking, for example, $\alpha = \frac{3}{2}$, $\beta = -\frac{1}{2}$, we have $\lambda_1 = 1$, $\lambda_2 = 2$. Consider the equation
        \[
            \mathcal F(\vec x) = (x_1')^2 + 2(x_2')^2 = 1
        \]
        -- this defines an ellipse along the principal axes.
        DIAGRAM: graph
        Contrarily, if $\alpha = -\frac{1}{2}$, $\beta = \frac{3}{2}$, $\lambda_1 = 1$ and $\lambda_2 = -2$. Then
        \[
            \mathcal F(\vec x) = (x_1')^2 - 2(x_2')^2 = 1
        \]
        defines a hyperbola, again along the principal axes.
        DIAGRAM: graph
    \end{example}
    
    \lecture{Lecture 22 - 27/11/25}

    Let's consider 3 dimensions, so
    \[
        \mathcal F(\vec x) = \vec x^\intercal A \vec x = \lambda_1(\vec x_1')^2 + \lambda_2(\vec x_2')^2 + \lambda_3(\vec x_3')^2
    \]
    If $\lambda_1$, $\lambda_2$, $\lambda_3 > 0$, then $\mathcal F(\vec x) = 1$ defines an \emph{ellipsoid}.

    \begin{example}
        Consider the quadratic form $\mathcal F$ defined by
        \[
            A = \begin{pmatrix}
             0 & 1 & 0 \\
             1 & 0 & 1 \\
             1 & 1 & 0 \\
            \end{pmatrix}
        \]
        with eigenvalues $\lambda_1 = \lambda_2 = -1$ and $\lambda_3 = 2$, and eigenvectors 
        \[
            \vec u_1 = \frac{1}{\sqrt 2} \begin{pmatrix}
                1 \\
                -1 \\
                0
            \end{pmatrix}, \; \vec u_2 = \frac{1}{\sqrt 6} \begin{pmatrix}
                1 \\
                1 \\
                -2
            \end{pmatrix}, \; \vec u_3 = \frac{1}{\sqrt 3} \begin{pmatrix}
                1 \\
                1 \\
                1
            \end{pmatrix}
        \]
        Then
        \begin{align*}
            \mathcal F(\vec x) &= 2x_1x_2 + 2x_2x_3 + 2x_3x_1 \\
            &= -(x_1')^2 - (x_2')^2 + 2(x_3')^2
        \end{align*}
        and here $\mathcal F = 1$ defines a \emph{hyperboloid} (2-sheeted).
        DIAGRAM: graph, minimum and maximum of sheets labelled

        On the other hand, $\mathcal F = -1$ defines a 1-sheeted hyperboloid.
        DIAGRAM: circle radius 1 at origin.
    \end{example}
    
    \begin{remark}
        Why does the defining matrix of a quadratic form have to be symmetric? Write $M = S + A$ where $S$ and $A$ are the symmetric and antisymmetric parts respectively. Then one can see $\vec x^\intercal A \vec x = 0$, so
        \[
            \vec x^\intercal M \vec x = \vec x^\intercal (S + A) \vec x = \vec x^\intercal S \vec x
        \]
        This is essentially a manifestation of the commutativity of multiplication.
    \end{remark}
    
    \subsection{Quadrics and conics}

    We aim to classify solutions to quadrics and conics -- in particular, we classify up to geometric equivalence, meaning we make no distinction between isometries in $\RR^n$.

    \subsubsection{Quadrics}

    \begin{definition}
        A \emph{quadric} in $\RR^n$ is a hypersurface defined by
        \[
            Q(\vec x) = \vec x^\intercal A \vec x + \vec b^\intercal \vec x + c = 0
        \]
        for a real symmetric matrix $A$, a vector $b \in \RR^n$, and a scalar $c \in \RR$. Then
        \[
            Q(\vec x) = A_{ij}x_ix_j + b_ix_i + c = 0
        \]
    \end{definition}
    If $A$ is invertible, then we can complete the square. Take
    \[
        \vec y = \vec x + \frac 1 2 A^{-1} \vec b
    \]
    so $\vec y^\intercal = \vec x^\intercal + \frac 1 2 A^{-1} \vec b^\intercal$ -- note that $(A^{-1})^\intercal = A^{-1}$ since $A$ symmetric $\implies A^{-1}$ symmetric. Then
    \begin{align*}
        \vec y^\intercal A \vec y &= \vec x^\intercal A \vec x + \vec b^\intercal \vec x + \frac{1}{4} \vec b^\intercal A^{-1} \vec b \\
        &= Q(\vec x) + \underbrace{\left(\frac{1}{4} \vec b^\intercal A^{-1} \vec b - c\right)}_k
    \end{align*}
    Letting $\mathcal F$ be the quadratic form defined by $A$, we have
    \[
        \mathcal F(\vec y) = Q(\vec x) + k 
    \]
    and so
    \[
        Q(\vec x) = 0 \iff \mathcal F(\vec y) = k
    \]
    Now, diagonalising $\mathcal F$ gives us the principal axes as before, and the geometrical nature of the quadric is determined by the eigenvalues of $A$ and the value of $k$.

    \begin{itemize}
        \item All eigenvalues $>0$ and $k>0$: \emph{ellipsoid}.
        \item Eigenvalues of both signs, $k \neq 0$: \emph{hyperboloid}
    \end{itemize}
    
    If $A$ is not invertible or, equivalently, has a zero eigenvalue, we must analyse the quadric directly.

    \subsubsection{Conics}

    \begin{definition}
        A conic is a quadric in $\RR^2$.
    \end{definition}
    
    \begin{itemize}
        \item If $A$ is invertible, we end up with
        \[
            \lambda_1 (x_1')^2 + \lambda_2 (x_2')^2 = k
        \]
        by the analysis we've just done. Then
        \begin{itemize}
            \item If $\lambda_1$, $\lambda_2 > 0$, then $k > 0$, $k = 0$ and $k < 0$ yield an ellipse, a point, and no solution, respectively.
            \item If $\lambda_1$ and $\lambda_2$ are of opposite sign, then $k \neq 0$ yields a hyperbola, and $k = 0$ yields a pair of lines.
        \end{itemize}
        \item If $A$ is not invertible, things are more complicated. Assume $\lambda_1 > 0$, $\lambda_2 = 0$. Since $A$ is symmetric, we can still diagonalise, so $\vec x^\intercal A \vec x = \lambda_1(x_1')^2$. Thus our quadric becomes
        \begin{align*}
            \lambda_1(x_1')^2 + b_1'x_1' + b_2'x_2' + c &= 0 \\
            \implies \lambda_1 \underbrace{\left(x_1' + \frac{1}{2 \lambda_1}b_1'\right)}_{x_1''} {^2} + b_2x_2' + \underbrace{c - \frac{1}{4\lambda_1}(b_1')^2}_{c'} &= 0
        \end{align*}    
        Now, we have some more cases:
        \begin{itemize}
            \item If $b_2'=0$, then $c' < 0$, $c = 0$ and $c > 0$ yield a pair of lines, a line and no solutions, respectively.
            \item If $b_2'\neq 0$, we have
            \[
                \lambda_1(x_1'')^2 + b_2'x_2''
            \]
            for $x_2'' = x_2' + \frac{c'}{b_2'}$. This is a parabola.
        \end{itemize}
    \end{itemize}
    
    Importantly, all the transformations we have done preserve distances and angles, so the original conic's solutions are geometrically equivalent.

    Now, we wish to look at conics in Cartesian coordinates. Let us consider some examples.

    \begin{itemize}
        \item \emph{Ellipses}. Consider
        \[
            \frac{x^2}{a^2} + \frac{y^2}{b^2} = 1
        \]
        Supposing $a > b$, $a$ is the \emph{semi-major axis} and $b$ is the \emph{semi-minor axis}. The \emph{eccentricity} $0 \leq e < 1$ satisfies $b^2 = a^2(1-e^2)$, and the \emph{foci} are at $x = \pm ae$. 

        DIAGRAM: graph, foci labelled
        \item \emph{Parabolas}. Consider
        \[
            y^2 = 4ax
        \]
        A parabola has eccentricity 1, and focus at $x=a$.
        \item \emph{Hyperbolas}. Consider
        \[
            \frac{x^2}{a^2} - \frac{y^2}{b^2} = \pm 1
        \]
        Now the eccentricity $1 < e$ satisfies $b^2 = a^2(e^2-1)$, and the foci are again at $\pm ae$.
    \end{itemize}
    
    \lecture{Lecture 23 -- 29/11/25}

    \lecture{Lecture 24 -- 02/12/25}

    For a rotation matrix $R$, consider a vector $\vec x$ which is taken to $\vec x'$ satisfying $x'_i = R_{ij}x_j$. Then, we have two ways to view this transformation:
    
    \begin{itemize}
        \item Transformation of vectors (active point of view). 
        DIAGRAM: $\vec e_1$, $\vec e_2$ axes, $\vec x$, $\vec x'$, $|\vec x| = |\vec x'|$, $x_i'$ components of new vector with respect to standard basis.
        \item Change of basis (passive point of view).
        DIAGRAM: $\vec e_1$, $\vec e_2$ axes, rotate to another set of axes $\vec u_1$, $\vec u_2$, $\vec x$ remains unchanged, but coordinates of $\vec x$ change (coordinates labelled in both basis systems). $x_i'$ components of the same vector $\vec x$ but with respect to the new basis spanned by $\vec u_1$, $\vec u_2$. Furthermore, we have
        \[
            \vec u_i = \sum_j (R^{-1})_{ji} \vec e_j = \sum_j R_{ij} \vec e_j
        \]
    \end{itemize}
    
    \begin{remark*}
        Comparing this to the standard change of basis approach, we have $P = R^{-1}$.
    \end{remark*}
    
    \subsubsection{2D Minkowski space and Lorentz transformations}

    Consider the `inner product' on $\RR^2$ given by
    \[
        (\vec x, \vec y) = \vec x^\intercal J \vec y
    \]
    where
    \[
        J = \begin{pmatrix}
         1 & 0 \\
         0 & -1 \\
        \end{pmatrix}
    \]
    so if $\vec x = (x_0, x_1)^\intercal$ and $\vec y = (y_0, y_1)^\intercal$, $(\vec x, \vec y) = x_0y_0 - x_1y_1$. Note that this is not an inner product since it is not positive-definite (but it is still symmetric and bilinear).

    Consider now the basis
    \[
        \vec e_0 = \begin{pmatrix}
         1 \\
         0 \\
        \end{pmatrix}, \; \vec e_1 = \begin{pmatrix}
         0 \\
         1 \\
        \end{pmatrix}
    \]
    which is `orthonormal' in the sense that $(\vec e_0, \vec e_0) = 1$, $-(\vec e_1, \vec e_1) = 1$, $(\vec e_0, \vec e_1) = 0$.

    \begin{definition}[Minkowski space]
        This inner product is called the \emph{Minkowski metric}, and $\RR^2$ equipped with the Minkowski metric is called \emph{Minkowski space}.
    \end{definition}
    Consider a matrix of the form
    \[
        M = \begin{pmatrix}
         M_{00} & M_{01} \\
         M_{10} & M_{11} \\
        \end{pmatrix}
    \]
    associated to a linear map $T: \RR^2 \to \RR^2$. When does $M$ preserve the Minkowski metric? We need $(M \vec x, M \vec y) = (\vec x, \vec y)$ for all $\vec x$, $\vec y \in \RR^2$. So
    \begin{align*}
        \vec x^\intercal J \vec y &= (M \vec x)^\intercal J (M \vec y) \\
        &= \vec x^\intercal (M^\intercal J M)\vec y 
    \end{align*}
    so we need $M^\intercal J M = J$. The matrices $M$ for which this holds form a group. Furthermore, we can see from the above that $\det M = \pm 1$.
    
    \begin{definition}[Lorentz group]
        The \emph{Lorentz group} is the subgroup of the above group satisfying $\det M = 1$ and $M_{00} > 0$.
    \end{definition}
    Let us now determine the general form of $M$ in the Lorentz group, in much the same method as for orthogonal matrices.
    \begin{itemize}
        \item $(\vec e_0, \vec e_0) = 1$, and so we need $(M \vec e_0, M \vec e_0) = 1$. We need
        \begin{align*}
            1 &= \begin{pmatrix}
             1 & 0
            \end{pmatrix} \begin{pmatrix}
             M_{00} & M_{01} \\
             M_{10} & M_{11} \\
            \end{pmatrix} \begin{pmatrix}
             1 & 0 \\
             0 & -1 \\
            \end{pmatrix} \begin{pmatrix}
             M_{00} & M_{01} \\
             M_{10} & M_{11} \\
            \end{pmatrix} \begin{pmatrix}
             1 \\
             0 \\
            \end{pmatrix} \\
            &= M_{00}^2 - M_{10}^2
        \end{align*}
        \item Doing the same thing for $\vec e_0$, we get $M_{01}^2 - M_{11}^2 = -1$.
        \item Doing the same thing for $\vec e_0$ and $\vec e_1$, we have $M_{00}M_{01}-M_{10}M_{11}=0$.
    \end{itemize}
    Combining all of these results, we have
    \[
        M = \begin{pmatrix}
         \cosh \theta & \sinh \theta \\
         \sinh \theta & \cosh \theta \\
        \end{pmatrix}
    \]
    for some $\theta \in \RR$. We have $M(\theta_1)M(\theta_2) = M(\theta_1 + \theta_2)$ using hyperbolic angle sum formulae. What is the physical interpretation of Lorentz transformations?
    
    DIAGRAM: y axis x0, x axis x1. Vector $\vec x$ goes to $\vec x'$.
    Upper/lower hyperbola: $k > 0$. $y=x$ or $y=-x$: $k = 0$. Right/left hyperbola $k < 0$.

    Fix $(\vec x, \vec x) = k$ constant -- we obtain the figures as shown above. Importantly, since Lorentz transformations preserve the `norm', they preserve the curve that $\vec x$ is on. The fact that $M_{00} > 0$ means that $\vec x$ and $\vec x'$ must in particular lie on the same branch of the hyperbola.

    Observe
    \begin{align*}
        M(\theta) &= \begin{pmatrix}
         \cosh \theta & \sinh \theta \\
         \sinh \theta & \cosh \theta \\
        \end{pmatrix} \\
        &= \cosh \theta \begin{pmatrix}
         1 & \tanh \theta \\
         \tanh \theta & 1 \\
        \end{pmatrix} \\
        &= \frac{1}{\sqrt{1-\tanh^2\theta}} \begin{pmatrix}
         1 & \tanh \theta \\
         \tanh \theta & 1 \\
        \end{pmatrix}
    \end{align*}
    and we put $v = \tanh \theta$, with the idea that $v$ represents velocity, taking the speed of light $c=1$. Then $-1 < v < 1$, which makes sense physically. Rename $x_0 \to t$, $x_1 \to x$, the time and space coordinates respectively.
    
    Then
    \[
        \begin{cases}
            t' = \frac{1}{\sqrt{1-v^2}}(t+vx) \\
            x' = \frac{1}{\sqrt{1-v^2}}(x+vt)
        \end{cases}
    \]
    so Lorentz transformations modify the time and space coordinates for observers moving with relative velocity $v$ in special relativity. Then, the passive point of view corresponds to changing one's frame of reference!

    The factor $\frac{1}{\sqrt{1-v^2}}$ in Lorentz transformations is called the \emph{Lorentz factor} $\gamma$ and gives rise to time dilation and length contraction effects.

    Also, we see that composition of velocity works differently: since $M(\theta_1)M(\theta_2) = M(\theta_3)$ where $\theta_3 = \theta_1 + \theta_2$, we have, recalling that $v_i = \tanh \theta_i$,
    \[
        v_3 = \frac{v_1+v_2}{1+v_1v_2}
    \]
    and this is how we add velocities in special relativity (to ensure consistency with $|v_i| < 1$).

    This concludes the course. $\blacksquare$
    
    

\end{document}
